% Options for packages loaded elsewhere
% Options for packages loaded elsewhere
\PassOptionsToPackage{unicode}{hyperref}
\PassOptionsToPackage{hyphens}{url}
\PassOptionsToPackage{dvipsnames,svgnames,x11names}{xcolor}
%
\documentclass[
  brazilian,
  a4paper,
  oneside]{scrbook}
\usepackage{xcolor}
\usepackage{amsmath,amssymb}
\setcounter{secnumdepth}{5}
\usepackage{iftex}
\ifPDFTeX
  \usepackage[T1]{fontenc}
  \usepackage[utf8]{inputenc}
  \usepackage{textcomp} % provide euro and other symbols
\else % if luatex or xetex
  \usepackage{unicode-math} % this also loads fontspec
  \defaultfontfeatures{Scale=MatchLowercase}
  \defaultfontfeatures[\rmfamily]{Ligatures=TeX,Scale=1}
\fi
\usepackage{lmodern}
\ifPDFTeX\else
  % xetex/luatex font selection
\fi
% Use upquote if available, for straight quotes in verbatim environments
\IfFileExists{upquote.sty}{\usepackage{upquote}}{}
\IfFileExists{microtype.sty}{% use microtype if available
  \usepackage[]{microtype}
  \UseMicrotypeSet[protrusion]{basicmath} % disable protrusion for tt fonts
}{}
\makeatletter
\@ifundefined{KOMAClassName}{% if non-KOMA class
  \IfFileExists{parskip.sty}{%
    \usepackage{parskip}
  }{% else
    \setlength{\parindent}{0pt}
    \setlength{\parskip}{6pt plus 2pt minus 1pt}}
}{% if KOMA class
  \KOMAoptions{parskip=half}}
\makeatother
% Make \paragraph and \subparagraph free-standing
\makeatletter
\ifx\paragraph\undefined\else
  \let\oldparagraph\paragraph
  \renewcommand{\paragraph}{
    \@ifstar
      \xxxParagraphStar
      \xxxParagraphNoStar
  }
  \newcommand{\xxxParagraphStar}[1]{\oldparagraph*{#1}\mbox{}}
  \newcommand{\xxxParagraphNoStar}[1]{\oldparagraph{#1}\mbox{}}
\fi
\ifx\subparagraph\undefined\else
  \let\oldsubparagraph\subparagraph
  \renewcommand{\subparagraph}{
    \@ifstar
      \xxxSubParagraphStar
      \xxxSubParagraphNoStar
  }
  \newcommand{\xxxSubParagraphStar}[1]{\oldsubparagraph*{#1}\mbox{}}
  \newcommand{\xxxSubParagraphNoStar}[1]{\oldsubparagraph{#1}\mbox{}}
\fi
\makeatother

\usepackage{color}
\usepackage{fancyvrb}
\newcommand{\VerbBar}{|}
\newcommand{\VERB}{\Verb[commandchars=\\\{\}]}
\DefineVerbatimEnvironment{Highlighting}{Verbatim}{commandchars=\\\{\}}
% Add ',fontsize=\small' for more characters per line
\usepackage{framed}
\definecolor{shadecolor}{RGB}{241,243,245}
\newenvironment{Shaded}{\begin{snugshade}}{\end{snugshade}}
\newcommand{\AlertTok}[1]{\textcolor[rgb]{0.68,0.00,0.00}{#1}}
\newcommand{\AnnotationTok}[1]{\textcolor[rgb]{0.37,0.37,0.37}{#1}}
\newcommand{\AttributeTok}[1]{\textcolor[rgb]{0.40,0.45,0.13}{#1}}
\newcommand{\BaseNTok}[1]{\textcolor[rgb]{0.68,0.00,0.00}{#1}}
\newcommand{\BuiltInTok}[1]{\textcolor[rgb]{0.00,0.23,0.31}{#1}}
\newcommand{\CharTok}[1]{\textcolor[rgb]{0.13,0.47,0.30}{#1}}
\newcommand{\CommentTok}[1]{\textcolor[rgb]{0.37,0.37,0.37}{#1}}
\newcommand{\CommentVarTok}[1]{\textcolor[rgb]{0.37,0.37,0.37}{\textit{#1}}}
\newcommand{\ConstantTok}[1]{\textcolor[rgb]{0.56,0.35,0.01}{#1}}
\newcommand{\ControlFlowTok}[1]{\textcolor[rgb]{0.00,0.23,0.31}{\textbf{#1}}}
\newcommand{\DataTypeTok}[1]{\textcolor[rgb]{0.68,0.00,0.00}{#1}}
\newcommand{\DecValTok}[1]{\textcolor[rgb]{0.68,0.00,0.00}{#1}}
\newcommand{\DocumentationTok}[1]{\textcolor[rgb]{0.37,0.37,0.37}{\textit{#1}}}
\newcommand{\ErrorTok}[1]{\textcolor[rgb]{0.68,0.00,0.00}{#1}}
\newcommand{\ExtensionTok}[1]{\textcolor[rgb]{0.00,0.23,0.31}{#1}}
\newcommand{\FloatTok}[1]{\textcolor[rgb]{0.68,0.00,0.00}{#1}}
\newcommand{\FunctionTok}[1]{\textcolor[rgb]{0.28,0.35,0.67}{#1}}
\newcommand{\ImportTok}[1]{\textcolor[rgb]{0.00,0.46,0.62}{#1}}
\newcommand{\InformationTok}[1]{\textcolor[rgb]{0.37,0.37,0.37}{#1}}
\newcommand{\KeywordTok}[1]{\textcolor[rgb]{0.00,0.23,0.31}{\textbf{#1}}}
\newcommand{\NormalTok}[1]{\textcolor[rgb]{0.00,0.23,0.31}{#1}}
\newcommand{\OperatorTok}[1]{\textcolor[rgb]{0.37,0.37,0.37}{#1}}
\newcommand{\OtherTok}[1]{\textcolor[rgb]{0.00,0.23,0.31}{#1}}
\newcommand{\PreprocessorTok}[1]{\textcolor[rgb]{0.68,0.00,0.00}{#1}}
\newcommand{\RegionMarkerTok}[1]{\textcolor[rgb]{0.00,0.23,0.31}{#1}}
\newcommand{\SpecialCharTok}[1]{\textcolor[rgb]{0.37,0.37,0.37}{#1}}
\newcommand{\SpecialStringTok}[1]{\textcolor[rgb]{0.13,0.47,0.30}{#1}}
\newcommand{\StringTok}[1]{\textcolor[rgb]{0.13,0.47,0.30}{#1}}
\newcommand{\VariableTok}[1]{\textcolor[rgb]{0.07,0.07,0.07}{#1}}
\newcommand{\VerbatimStringTok}[1]{\textcolor[rgb]{0.13,0.47,0.30}{#1}}
\newcommand{\WarningTok}[1]{\textcolor[rgb]{0.37,0.37,0.37}{\textit{#1}}}

\usepackage{longtable,booktabs,array}
\newcounter{none} % for unnumbered tables
\usepackage{calc} % for calculating minipage widths
% Correct order of tables after \paragraph or \subparagraph
\usepackage{etoolbox}
\makeatletter
\patchcmd\longtable{\par}{\if@noskipsec\mbox{}\fi\par}{}{}
\makeatother
% Allow footnotes in longtable head/foot
\IfFileExists{footnotehyper.sty}{\usepackage{footnotehyper}}{\usepackage{footnote}}
\makesavenoteenv{longtable}
\usepackage{graphicx}
\makeatletter
\newsavebox\pandoc@box
\newcommand*\pandocbounded[1]{% scales image to fit in text height/width
  \sbox\pandoc@box{#1}%
  \Gscale@div\@tempa{\textheight}{\dimexpr\ht\pandoc@box+\dp\pandoc@box\relax}%
  \Gscale@div\@tempb{\linewidth}{\wd\pandoc@box}%
  \ifdim\@tempb\p@<\@tempa\p@\let\@tempa\@tempb\fi% select the smaller of both
  \ifdim\@tempa\p@<\p@\scalebox{\@tempa}{\usebox\pandoc@box}%
  \else\usebox{\pandoc@box}%
  \fi%
}
% Set default figure placement to htbp
\def\fps@figure{htbp}
\makeatother


% definitions for citeproc citations
\NewDocumentCommand\citeproctext{}{}
\NewDocumentCommand\citeproc{mm}{%
  \begingroup\def\citeproctext{#2}\cite{#1}\endgroup}
\makeatletter
 % allow citations to break across lines
 \let\@cite@ofmt\@firstofone
 % avoid brackets around text for \cite:
 \def\@biblabel#1{}
 \def\@cite#1#2{{#1\if@tempswa , #2\fi}}
\makeatother
\newlength{\cslhangindent}
\setlength{\cslhangindent}{1.5em}
\newlength{\csllabelwidth}
\setlength{\csllabelwidth}{3em}
\newenvironment{CSLReferences}[2] % #1 hanging-indent, #2 entry-spacing
 {\begin{list}{}{%
  \setlength{\itemindent}{0pt}
  \setlength{\leftmargin}{0pt}
  \setlength{\parsep}{0pt}
  % turn on hanging indent if param 1 is 1
  \ifodd #1
   \setlength{\leftmargin}{\cslhangindent}
   \setlength{\itemindent}{-1\cslhangindent}
  \fi
  % set entry spacing
  \setlength{\itemsep}{#2\baselineskip}}}
 {\end{list}}
\usepackage{calc}
\newcommand{\CSLBlock}[1]{\hfill\break\parbox[t]{\linewidth}{\strut\ignorespaces#1\strut}}
\newcommand{\CSLLeftMargin}[1]{\parbox[t]{\csllabelwidth}{\strut#1\strut}}
\newcommand{\CSLRightInline}[1]{\parbox[t]{\linewidth - \csllabelwidth}{\strut#1\strut}}
\newcommand{\CSLIndent}[1]{\hspace{\cslhangindent}#1}

\ifLuaTeX
\usepackage[bidi=basic,shorthands=off]{babel}
\else
\usepackage[bidi=default,shorthands=off]{babel}
\fi
\ifLuaTeX
  \usepackage{selnolig} % disable illegal ligatures
\fi


\setlength{\emergencystretch}{3em} % prevent overfull lines

\providecommand{\tightlist}{%
  \setlength{\itemsep}{0pt}\setlength{\parskip}{0pt}}



 


\usepackage[brazil]{babel}
\makeatletter
\@ifpackageloaded{tcolorbox}{}{\usepackage[skins,breakable]{tcolorbox}}
\@ifpackageloaded{fontawesome5}{}{\usepackage{fontawesome5}}
\definecolor{quarto-callout-color}{HTML}{909090}
\definecolor{quarto-callout-note-color}{HTML}{0758E5}
\definecolor{quarto-callout-important-color}{HTML}{CC1914}
\definecolor{quarto-callout-warning-color}{HTML}{EB9113}
\definecolor{quarto-callout-tip-color}{HTML}{00A047}
\definecolor{quarto-callout-caution-color}{HTML}{FC5300}
\definecolor{quarto-callout-color-frame}{HTML}{acacac}
\definecolor{quarto-callout-note-color-frame}{HTML}{4582ec}
\definecolor{quarto-callout-important-color-frame}{HTML}{d9534f}
\definecolor{quarto-callout-warning-color-frame}{HTML}{f0ad4e}
\definecolor{quarto-callout-tip-color-frame}{HTML}{02b875}
\definecolor{quarto-callout-caution-color-frame}{HTML}{fd7e14}
\makeatother
\makeatletter
\@ifpackageloaded{bookmark}{}{\usepackage{bookmark}}
\makeatother
\makeatletter
\@ifpackageloaded{caption}{}{\usepackage{caption}}
\AtBeginDocument{%
\ifdefined\contentsname
  \renewcommand*\contentsname{Índice}
\else
  \newcommand\contentsname{Índice}
\fi
\ifdefined\listfigurename
  \renewcommand*\listfigurename{Lista de Figuras}
\else
  \newcommand\listfigurename{Lista de Figuras}
\fi
\ifdefined\listtablename
  \renewcommand*\listtablename{Lista de Tabelas}
\else
  \newcommand\listtablename{Lista de Tabelas}
\fi
\ifdefined\figurename
  \renewcommand*\figurename{Figura}
\else
  \newcommand\figurename{Figura}
\fi
\ifdefined\tablename
  \renewcommand*\tablename{Tabela}
\else
  \newcommand\tablename{Tabela}
\fi
}
\@ifpackageloaded{float}{}{\usepackage{float}}
\floatstyle{ruled}
\@ifundefined{c@chapter}{\newfloat{codelisting}{h}{lop}}{\newfloat{codelisting}{h}{lop}[chapter]}
\floatname{codelisting}{Listagem}
\newcommand*\listoflistings{\listof{codelisting}{Lista de Listagens}}
\makeatother
\makeatletter
\makeatother
\makeatletter
\@ifpackageloaded{caption}{}{\usepackage{caption}}
\@ifpackageloaded{subcaption}{}{\usepackage{subcaption}}
\makeatother
\usepackage{bookmark}
\IfFileExists{xurl.sty}{\usepackage{xurl}}{} % add URL line breaks if available
\urlstyle{same}
\hypersetup{
  pdftitle={RHCgen},
  pdfauthor={Anderson Lineu Siqueira dos Santos; Tatyellen Natasha da Costa Oliveira; Raphael de Freitas Saldanha; Raquel de Vasconcellos Carvalhaes de Oliveira},
  pdflang={pt-BR},
  colorlinks=true,
  linkcolor={Maroon},
  filecolor={Maroon},
  citecolor={Blue},
  urlcolor={Blue},
  pdfcreator={LaTeX via pandoc}}


\title{RHCgen}
\usepackage{etoolbox}
\makeatletter
\providecommand{\subtitle}[1]{% add subtitle to \maketitle
  \apptocmd{\@title}{\par {\large #1 \par}}{}{}
}
\makeatother
\subtitle{Um Pacote R para Automação, Padronização e Avaliação da
Qualidade dos Dados do Registro Hospitalar de Câncer no Brasil}
\author{Anderson Lineu Siqueira dos Santos \and Tatyellen Natasha da
Costa Oliveira \and Raphael de Freitas Saldanha \and Raquel de
Vasconcellos Carvalhaes de Oliveira}
\date{01/05/2026}
\begin{document}
\frontmatter
\maketitle

\renewcommand*\contentsname{Índice}
{
\setcounter{tocdepth}{2}
\tableofcontents
}

\mainmatter
\bookmarksetup{startatroot}

\chapter*{Bem-vindo}\label{bem-vindo}
\addcontentsline{toc}{chapter}{Bem-vindo}

\markboth{Bem-vindo}{Bem-vindo}

Esta nota técnica descreve o \textbf{RHCgen}, um pacote computacional de
acesso aberto implementado na linguagem R, voltado à automação,
padronização e avaliação da qualidade dos dados do \textbf{Registro
Hospitalar de Câncer (RHC)} no Brasil.

O pacote está disponível no GitHub:
\href{https://github.com/andersonlineu/RHCgen}{github.com/andersonlineu/RHCgen}

Um guia de uso está disponível em:
\href{https://andersonlineu.github.io/GuiaRHCgen/}{andersonlineu.github.io/GuiaRHCgen}

\begin{center}\rule{0.5\linewidth}{0.5pt}\end{center}

\section*{Como citar}\label{como-citar}
\addcontentsline{toc}{section}{Como citar}

\markright{Como citar}

\begin{quote}
Santos, A. L. S. dos., Oliveira, T. N. da C., Saldanha, R. de F., \&
Oliveira, R. de V. C. de. (2026). \emph{RHCgen: Um Pacote R para
Automação, Padronização e Avaliação da Qualidade dos Dados do Registro
Hospitalar de Câncer no Brasil.} Zenodo.
\url{https://doi.org/10.5281/zenodo.19962293}
\end{quote}

\subsection*{BibTeX}\label{bibtex}
\addcontentsline{toc}{subsection}{BibTeX}

\begin{Shaded}
\begin{Highlighting}[]
\VariableTok{@techreport}\NormalTok{\{}\OtherTok{santos\_rhcgen\_2026}\NormalTok{,}
  \DataTypeTok{author}\NormalTok{    = \{Santos, Anderson Lineu Siqueira dos and}
\NormalTok{               Oliveira, Tatyellen Natasha da Costa and}
\NormalTok{               Saldanha, Raphael de Freitas and}
\NormalTok{               Oliveira, Raquel de Vasconcellos Carvalhaes de\},}
  \DataTypeTok{title}\NormalTok{     = \{\{RHCgen\}: Um Pacote \{R\} para Automação, Padronização e}
\NormalTok{               Avaliação da Qualidade dos Dados do Registro Hospitalar}
\NormalTok{               de Câncer no Brasil\},}
  \DataTypeTok{year}\NormalTok{      = \{2026\},}
  \DataTypeTok{doi}\NormalTok{       = \{10.5281/zenodo.19962293\},}
  \DataTypeTok{publisher}\NormalTok{ = \{Zenodo\}}
\NormalTok{\}}
\end{Highlighting}
\end{Shaded}

\begin{center}\rule{0.5\linewidth}{0.5pt}\end{center}

\section*{Resumo}\label{resumo}
\addcontentsline{toc}{section}{Resumo}

\markright{Resumo}

\textbf{Introdução:} O Registro Hospitalar de Câncer (RHC) reúne dados
clínicos e assistenciais de pacientes oncológicos atendidos em hospitais
habilitados pelo Sistema Único de Saúde no Brasil, constituindo uma base
estratégica para a vigilância e a pesquisa em oncologia. Contudo, a
ausência de ferramentas padronizadas para seu processamento, aliada à
heterogeneidade estrutural e semântica dos arquivos disponibilizados,
impõe barreiras significativas à sua utilização em larga escala e de
forma reprodutível.

\textbf{Objetivo:} Descrever um pacote computacional de acesso aberto,
implementado na linguagem R, voltado à automação, padronização do
processamento e avaliação da qualidade dos dados do Registro Hospitalar
de Câncer no Brasil.

\textbf{Métodos:} Estudo metodológico que descreve o desenvolvimento do
pacote RHCgen, implementado em R com o auxílio dos pacotes
\emph{usethis}, \emph{devtools} e \emph{roxygen2}. O pacote foi
estruturado como um \emph{pipeline} automatizado de extração,
transformação e carga (ETL), composto por funções modulares para leitura
e consolidação de arquivos no formato \texttt{.dbf}, padronização
estrutural, variações semânticas e recodificação de variáveis, além de
avaliação de completude e consistência dos dados.

\textbf{Resultados:} O processamento de 24 arquivos nacionais do RHC
(2000--2023) resultou em uma base analítica consolidada com
\textbf{5.394.771 registros} e \textbf{59 variáveis}, com tempo de
execução de 4,61 minutos e consumo aproximado de 2,01 GB de memória RAM.

\textbf{Conclusão:} O RHCgen oferece uma solução padronizada,
reprodutível e de acesso aberto para o processamento de dados do RHC,
contribuindo para ampliar a produção científica baseada no RHC e
fortalecer a vigilância oncológica no Brasil.

\textbf{Palavras-chave:} Registros Hospitalares de Câncer; Processamento
de Dados; Reprodutibilidade dos Resultados; Sistemas de Informação em
Saúde; Neoplasias; Linguagem R; RHCgen.

\begin{center}\rule{0.5\linewidth}{0.5pt}\end{center}

\section*{Sobre os autores}\label{sobre-os-autores}
\addcontentsline{toc}{section}{Sobre os autores}

\markright{Sobre os autores}

\textbf{Anderson Lineu Siqueira dos Santos} --- Escola Nacional de Saúde
Pública Sergio Arouca, Fundação Oswaldo Cruz, Rio de Janeiro, Brasil.

\href{https://orcid.org/0000-0002-1703-9310}{0000-0002-1703-9310}

\textbf{Tatyellen Natasha da Costa Oliveira} --- Escola Nacional de
Saúde Pública Sergio Arouca, Fundação Oswaldo Cruz; Instituto Evandro
Chagas, Ananindeua, PA, Brasil. 
\href{https://orcid.org/0000-0002-4781-3952}{0000-0002-4781-3952}

\textbf{Raphael de Freitas Saldanha} --- Instituto de Comunicação e
Informação Científica e Tecnológica em Saúde, Fundação Oswaldo Cruz, Rio
de Janeiro, RJ, Brasil. 
\href{https://orcid.org/0000-0003-0652-8466}{0000-0003-0652-8466}

\textbf{Raquel de Vasconcellos Carvalhaes de Oliveira} --- Instituto
Nacional de Infectologia, Fundação Oswaldo Cruz, Rio de Janeiro, RJ,
Brasil. 
\href{https://orcid.org/0000-0001-9387-8645}{0000-0001-9387-8645}

\bookmarksetup{startatroot}

\chapter{Introdução}\label{introduuxe7uxe3o}

Nos últimos anos, observa-se aumento da quantidade de pesquisas que
envolvem bases de dados de câncer, a fim de melhor compreender a
dinâmica da evolução, tratamento e prevenção da doença ({Bray et al.}
2024). Iniciativas como o \emph{Global Cancer Observatory}, o
\emph{Cancer Incidence in Five Continents} e o \emph{CANCER TODAY}
possuem abrangência global e disponibilizam dados sobre incidência,
mortalidade e prevalência de câncer em 185 países e mais de 300 regiões
no mundo (SEER 2025; Cancer Today 2025). Em níveis locais, bases de
dados norte-americanas como o \emph{Surveillance, Epidemiology, and End
Results} do \emph{National Cancer Institute} e o \emph{National Cancer
Database} reúnem dados de mais de 1.500 hospitais, cobrindo cerca de
48\% da população dos EUA (SEER 2025; ACS 2025). Na Europa, o
\emph{European Cancer Information System}, operado pela Comissão
Europeia, coleta dados abrangentes de incidência, mortalidade e
sobrevivência do câncer, permitindo análises comparativas entre regiões
e subsidiando políticas de saúde pública eficazes (EORTC 2017; ECIS
2025).

O Brasil dispõe de duas bases de dados específicas para o registro de
informações sobre câncer: os \textbf{Registros de Câncer de Base
Populacional (RCBP)} e os \textbf{Registros Hospitalares de Câncer
(RHC)}, ambos mantidos pelo Instituto Nacional do Câncer (INCA).
Enquanto o RCBP registra a ocorrência e as características dos casos
novos na população, o RHC reúne informações clínicas e terapêuticas de
pacientes atendidos em hospitais habilitados na Atenção Especializada em
Oncologia do Sistema Único de Saúde (SUS). Os RHC são responsáveis pela
coleta, armazenamento, processamento, análise e divulgação contínua
desses dados, incluindo informações sobre identificação, características
demográficas, diagnóstico, tratamento e evolução dos tumores.
Obrigatórios em hospitais que atendem pacientes com câncer no SUS, esses
registros subsidiam a gestão hospitalar, o acompanhamento clínico e a
avaliação da qualidade da assistência e dos desfechos dos tratamentos
(INCA 2000).

No contexto da análise de dados em saúde no Brasil, destaca-se a
ferramenta \textbf{TabWin}, desenvolvida pelo DATASUS (Brasil 2025), que
desempenha um papel fundamental ao permitir a tabulação e o tratamento
desses dados. Entretanto, apesar da sua capacidade de leitura de
arquivos no formato DBC, construção de tabelas, indicadores, gráficos e
mapas, o TabWin é restrito ao sistema operacional \emph{Windows} e tem
capacidades analíticas limitadas. Para estatísticas mais complexas
existe a necessidade de exportação dos dados para pacotes estatísticos
mais robustos, o que adiciona trabalho e potencial para erros (Jorge et
al. 2010; Silva 2009).

Essas limitações têm impulsionado a busca por ferramentas mais
eficientes e reprodutíveis para análise de dados, viabilizada por
ferramentas de ciência de dados que têm transformado a preparação dos
dados em pesquisa. A \textbf{linguagem R}, amplamente utilizada para
análise estatística, destaca-se por permitir a organização e manipulação
de grandes volumes de dados heterogêneos, comuns em epidemiologia em
saúde pública, ao facilitar a conversão de dados brutos em formatos
prontos para análise e contribuir para a produção de resultados
científicos confiáveis ({Lowndes et al.} 2017).

Alguns pacotes em R específicos para manipulação e processamento de
dados de saúde no Brasil facilitam o acesso e análise de informações de
saúde pública, como o \textbf{microdatasus} (Saldanha et al. 2019) ---
que viabiliza o download e pré-processamento de dados de alguns Sistemas
de Informação do SUS diretamente do DATASUS --- e o \textbf{Rtabnetsp}
(Morais et al. 2021), que permite a extração de indicadores de saúde dos
servidores TabNet da Secretaria de Estado da Saúde de São Paulo.
Entretanto, até o momento não há uma ferramenta que permita a automação
e padronização integral do processamento do RHC.

Como consequência, a utilização desses dados em larga escala permanece
limitada por barreiras relevantes de interoperabilidade e
reprodutibilidade. Na prática, a ausência de um \emph{pipeline}
estruturado obriga cada pesquisador a desenvolver rotinas próprias de
limpeza, transformação e harmonização --- geralmente baseadas em
\emph{scripts ad hoc}. Esse cenário fragmentado implica não apenas
repetição desnecessária das mesmas etapas de processamento por
diferentes pesquisadores, mas também a introdução de decisões analíticas
heterogêneas, que podem gerar estimativas divergentes a partir de uma
mesma base de dados.

Diante desse cenário, a proposta de um pacote que estabeleça um padrão
de governança dos dados torna-se particularmente relevante, ao garantir
que o processo de transformação do dado bruto em base analítica seja
\textbf{transparente, documentado e reprodutível}. A padronização dessas
etapas contribui para maior consistência entre estudos, além de
facilitar auditorias, validações externas e o compartilhamento de
métodos.

Esta nota técnica tem como objetivo descrever e apresentar o pacote
computacional de acesso aberto \textbf{RHCgen}, implementado na
linguagem R, voltado à automação, padronização do processamento e
avaliação da qualidade dos dados do Registro Hospitalar de Câncer no
Brasil, descrevendo sua arquitetura, funcionalidades e potencial de
aplicação em análises em larga escala com dados do RHC.

\bookmarksetup{startatroot}

\chapter{Fonte de Dados: o Registro Hospitalar de
Câncer}\label{fonte-de-dados-o-registro-hospitalar-de-cuxe2ncer}

\section{O que é o RHC?}\label{o-que-uxe9-o-rhc}

Os Registros Hospitalares de Câncer (RHC) em todo o território
brasileiro coletam dados individualizados de pacientes com diagnóstico
confirmado de câncer, obtidos de prontuários e outras fontes
hospitalares. Esses registros permitem tanto a análise detalhada em
nível individual quanto a avaliação agregada de desfechos clínicos e dos
padrões diagnósticos mais frequentes (INCA 2000).

\section{IntegradorRHC (IRHC)}\label{integradorrhc-irhc}

O RHCgen é específico para ser utilizado com dados do RHC
disponibilizados para \emph{download} no \textbf{IntegradorRHC (IRHC)},
no formato \emph{Data Base File} (\texttt{.dbf}), acessível em:

\begin{quote}
\url{https://irhc.inca.gov.br/RHCNet/}
\end{quote}

Implantado em \textbf{2007}, o IntegradorRHC é um sistema público que
consolida e disponibiliza informações processadas pelo Sistema de
Informação de Registro Hospitalar de Câncer (SisRHC), com base em
registros de pacientes atendidos em unidades habilitadas no SUS em todo
o território nacional.

A base de dados é \textbf{não nominativa} e estruturada a partir de
registros padronizados pelas instituições participantes, sendo o SisRHC
responsável pelos processos de entrada, validação e padronização das
informações (INCA 2011; Decit 2007).

\section{Cobertura temporal}\label{cobertura-temporal}

O RHCgen possui capacidade de processamento de arquivos provenientes de
diferentes períodos: \textbf{de 1985 a 2023}.

Para a construção do pacote e realização dos testes do \emph{pipeline},
foram utilizados dados provenientes do IntegradorRHC. O período
analisado compreendeu registros com ano do banco entre 1º de janeiro de
1985 e 31 de dezembro de 2023. Os dados foram extraídos em \textbf{11 de
julho de 2025}, considerando a versão da base disponível no sistema à
época, cuja última atualização havia ocorrido em 7 de abril de 2025,
assegurando a rastreabilidade dos dados utilizados.

\section{Dicionário de variáveis}\label{dicionuxe1rio-de-variuxe1veis}

O dicionário completo de variáveis do SisRHC está disponível no
IntegradorRHC:

\begin{quote}
\url{https://irhc.inca.gov.br/RHCNet/documentosTabWin.action?doc=dicionario.dados}
\end{quote}

\begin{tcolorbox}[enhanced jigsaw, arc=.35mm, bottomrule=.15mm, bottomtitle=1mm, breakable, colback=white, colbacktitle=quarto-callout-note-color!10!white, colframe=quarto-callout-note-color-frame, coltitle=black, left=2mm, leftrule=.75mm, opacityback=0, opacitybacktitle=0.6, rightrule=.15mm, title=\textcolor{quarto-callout-note-color}{\faInfo}\hspace{0.5em}{Rastreabilidade dos dados}, titlerule=0mm, toprule=.15mm, toptitle=1mm]

A identificação da data de extração e da versão da base são etapas
recomendadas para garantir a reprodutibilidade de análises conduzidas
com o RHCgen. O pacote cria automaticamente a variável
\texttt{Ano\_do\_Banco} a partir do nome original de cada arquivo
\texttt{.dbf}, permitindo identificar o ano de origem de cada registro.

\end{tcolorbox}

\bookmarksetup{startatroot}

\chapter{Desenvolvimento do RHCgen}\label{desenvolvimento-do-rhcgen}

\section{Plataforma e ferramentas}\label{plataforma-e-ferramentas}

O desenvolvimento do RHCgen foi realizado com o uso da \textbf{linguagem
de programação R} (versão 4.5.1) (R Core Team 2024), no ambiente de
desenvolvimento integrado \textbf{RStudio} (versão 2025.06.13) (RStudio
2024), com o auxílio dos seguintes pacotes:

{\def\LTcaptype{none} % do not increment counter
\begin{longtable}[]{@{}
  >{\raggedright\arraybackslash}p{(\linewidth - 2\tabcolsep) * \real{0.2353}}
  >{\raggedright\arraybackslash}p{(\linewidth - 2\tabcolsep) * \real{0.7647}}@{}}
\toprule\noalign{}
\begin{minipage}[b]{\linewidth}\raggedright
Pacote
\end{minipage} & \begin{minipage}[b]{\linewidth}\raggedright
Função no desenvolvimento
\end{minipage} \\
\midrule\noalign{}
\endhead
\bottomrule\noalign{}
\endlastfoot
\emph{usethis} & Configuração da estrutura do pacote; criação de
diretórios e arquivos padrão \\
\emph{devtools} & Construção, verificação e instalação do pacote; gestão
de dependências \\
\emph{roxygen2} & Geração automática da documentação das funções \\
\end{longtable}
}

\section{Disponibilidade e
instalação}\label{disponibilidade-e-instalauxe7uxe3o}

O RHCgen está disponível no GitHub com o objetivo de facilitar seu uso,
além de possibilitar a revisão e o aprimoramento por outros
pesquisadores:

\begin{quote}
\url{https://github.com/andersonlineu/RHCgen}
\end{quote}

Com o pacote \emph{remotes} previamente instalado no R, a instalação do
RHCgen pode ser realizada da seguinte forma:

\begin{Shaded}
\begin{Highlighting}[]
\NormalTok{remotes}\SpecialCharTok{::}\FunctionTok{install\_github}\NormalTok{(}\StringTok{"andersonlineu/RHCgen"}\NormalTok{)}
\end{Highlighting}
\end{Shaded}

\begin{tcolorbox}[enhanced jigsaw, arc=.35mm, bottomrule=.15mm, bottomtitle=1mm, breakable, colback=white, colbacktitle=quarto-callout-note-color!10!white, colframe=quarto-callout-note-color-frame, coltitle=black, left=2mm, leftrule=.75mm, opacityback=0, opacitybacktitle=0.6, rightrule=.15mm, title=\textcolor{quarto-callout-note-color}{\faInfo}\hspace{0.5em}{Dependência automática}, titlerule=0mm, toprule=.15mm, toptitle=1mm]

O pacote \emph{foreign} é instalado automaticamente no primeiro uso do
RHCgen. Esse pacote realiza a leitura de arquivos no formato
\texttt{.dbf}, converte campos de caracteres em fatores e respeita,
sempre que possível, os campos nulos.

\end{tcolorbox}

\section{Guia de uso}\label{guia-de-uso}

Um guia completo com orientações ao usuário sobre como utilizar o pacote
está disponível em:

\begin{quote}
\url{https://andersonlineu.github.io/GuiaRHCgen/}
\end{quote}

\section{Validação algorítmica e consistência dos
dados}\label{validauxe7uxe3o-algoruxedtmica-e-consistuxeancia-dos-dados}

A validação dos algoritmos de recodificação implementados no RHCgen foi
baseada na correspondência com as normas técnicas descritas em:

\begin{itemize}
\tightlist
\item
  \textbf{Manual} \emph{Registros Hospitalares de Câncer: Planejamento e
  Gestão} (INCA 2010)
\item
  \textbf{Dicionário de Variáveis do SisRHC} (INCA 2020)
\end{itemize}

Ambos publicados pelo Instituto Nacional de Câncer (INCA).

As rotinas de processamento foram construídas de forma
\textbf{determinística}, assegurando que a transformação dos dados
seguisse estritamente os padrões oficiais de codificação, incluindo
mapeamentos institucionais (CNES), classificações clínicas (CID-O) e
definições de desfechos.

A identificação dos estabelecimentos de saúde foi realizada por meio de
vinculação determinística utilizando o \textbf{código CNES} presente na
base do RHC como chave única para integração com a base oficial do CNES.

A confiabilidade das funções foi avaliada por meio de:

\begin{itemize}
\tightlist
\item
  \textbf{Testes unitários} --- valores de entrada previamente definidos
  foram comparados às saídas esperadas, com base nos dicionários
  oficiais
\item
  \textbf{Verificações de cobertura} --- para variáveis com grande
  número de categorias, garantindo que todos os códigos válidos do
  SisRHC fossem corretamente processados
\item
  \textbf{Testes de integração} --- desempenho das funções em execução
  sequencial ao longo do \emph{pipeline}
\item
  \textbf{Comparações com bases de referência} --- cotejo entre saídas
  automatizadas e bases revisadas manualmente
\end{itemize}

\begin{tcolorbox}[enhanced jigsaw, arc=.35mm, bottomrule=.15mm, bottomtitle=1mm, breakable, colback=white, colbacktitle=quarto-callout-important-color!10!white, colframe=quarto-callout-important-color-frame, coltitle=black, left=2mm, leftrule=.75mm, opacityback=0, opacitybacktitle=0.6, rightrule=.15mm, title=\textcolor{quarto-callout-important-color}{\faExclamation}\hspace{0.5em}{Dados reais na validação}, titlerule=0mm, toprule=.15mm, toptitle=1mm]

Diferente de abordagens baseadas em simulações, todas as etapas de
validação e estresse de volumetria utilizaram exclusivamente
\textbf{dados reais do RHC}, permitindo atestar a estabilidade e o
desempenho da ferramenta em condições operacionais reais.

\end{tcolorbox}

\bookmarksetup{startatroot}

\chapter{\texorpdfstring{O \emph{Pipeline} de
Processamento}{O Pipeline de Processamento}}\label{o-pipeline-de-processamento}

\section{\texorpdfstring{Processo único: a função
\texttt{construir\_banco()}}{Processo único: a função construir\_banco()}}\label{processo-uxfanico-a-funuxe7uxe3o-construir_banco}

A arquitetura do RHCgen é centrada na função modular
\texttt{construir\_banco()}, que opera como um \emph{pipeline}
automatizado de extração, transformação e carga --- \emph{automated
extract, transform, and load} (ETL).

Diferente de abordagens convencionais, esta função orquestra a transição
do dado bruto para a base analítica em \textbf{quatro estágios
críticos}:

\includegraphics[width=7.3in,height=7.9in]{04-pipeline_files/figure-latex/mermaid-figure-1.png}

\section{Os quatro estágios}\label{os-quatro-estuxe1gios}

\subsection{1. Harmonização
Estrutural}\label{harmonizauxe7uxe3o-estrutural}

Leitura e consolidação dinâmica de múltiplos arquivos \texttt{.dbf},
resolvendo inconsistências de nomenclatura e tipagem de colunas entre
diferentes séries históricas para garantir uma base única e consistente.

\subsection{2. Transformação de Tipos e
Atributos}\label{transformauxe7uxe3o-de-tipos-e-atributos}

Conversão automática de variáveis importadas erroneamente como fatores
para seus formatos originais (numérico para idade, data para eventos
temporais e caractere para os demais campos), assegurando a integridade
necessária para cálculos estatísticos.

\subsection{3. Enriquecimento e
Recodificação}\label{enriquecimento-e-recodificauxe7uxe3o}

Mapeamento de códigos técnicos --- como CID-O, CNES e Estadiamento
Clínico --- para termos descritivos completos, utilizando os dicionários
oficiais do INCA como referência (INCA 2010, 2020).

\subsection{4. Auditoria de Qualidade}\label{auditoria-de-qualidade}

Cálculo de indicadores de completude e consistência lógica para
avaliação da qualidade dos dados da base, permitindo identificar níveis
de preenchimento e incoerências entre variáveis, com base em escores
descritos por Romero e Cunha (2006) e Pinto et al. (2012).

\section{Uso básico}\label{uso-buxe1sico}

\begin{Shaded}
\begin{Highlighting}[]
\FunctionTok{library}\NormalTok{(RHCgen)}

\CommentTok{\# Certifique{-}se de que os arquivos .dbf estão no diretório de trabalho}
\CommentTok{\# com o padrão de nomenclatura rhcXX.dbf (ex: rhc20.dbf, rhc21.dbf)}

\NormalTok{dados\_RHC }\OtherTok{\textless{}{-}} \FunctionTok{construir\_banco}\NormalTok{()}
\end{Highlighting}
\end{Shaded}

\section{Relatório final de
processamento}\label{relatuxf3rio-final-de-processamento}

Ao término da execução, a função \texttt{construir\_banco()} exibe
automaticamente um relatório com os indicadores operacionais do
processamento:

{\def\LTcaptype{none} % do not increment counter
\begin{longtable}[]{@{}ll@{}}
\toprule\noalign{}
Indicador & Valor (exemplo --- base 2000--2023) \\
\midrule\noalign{}
\endhead
\bottomrule\noalign{}
\endlastfoot
Tempo de execução & 4,61 minutos \\
Memória utilizada & ≈ 2,01 GB \\
Total de registros & 5.394.771 \\
Total de variáveis & 59 \\
\end{longtable}
}

\begin{tcolorbox}[enhanced jigsaw, arc=.35mm, bottomrule=.15mm, bottomtitle=1mm, breakable, colback=white, colbacktitle=quarto-callout-tip-color!10!white, colframe=quarto-callout-tip-color-frame, coltitle=black, left=2mm, leftrule=.75mm, opacityback=0, opacitybacktitle=0.6, rightrule=.15mm, title=\textcolor{quarto-callout-tip-color}{\faLightbulb}\hspace{0.5em}{Produto final}, titlerule=0mm, toprule=.15mm, toptitle=1mm]

O dataframe resultante é mantido em memória e retornado ao ambiente de
execução via \texttt{return(data)}. Não há salvamento automático em
disco, o que oferece flexibilidade para que o usuário defina suas
próprias estratégias de exportação (por exemplo, \texttt{saveRDS()},
\texttt{write.csv()}, ou integração com bancos de dados).

\end{tcolorbox}

\bookmarksetup{startatroot}

\chapter{Leitura e Integração dos Arquivos
DBF}\label{leitura-e-integrauxe7uxe3o-dos-arquivos-dbf}

\section{\texorpdfstring{Função
\texttt{lerarquivoDBF()}}{Função lerarquivoDBF()}}\label{funuxe7uxe3o-lerarquivodbf}

\subsection{Finalidade}\label{finalidade}

A função \texttt{lerarquivoDBF()} realiza a \textbf{leitura automática}
dos arquivos DBF do Registro Hospitalar de Câncer disponíveis no
diretório de trabalho e combina esses arquivos em um único banco de
dados. Além disso, cria a variável \texttt{Ano\_do\_Banco}, extraída do
nome original do arquivo, permitindo identificar o ano de origem de cada
registro.

\subsection{Problema que resolve}\label{problema-que-resolve}

Os arquivos do RHC são disponibilizados separadamente, geralmente por
ano, em formato DBF. Esse formato exige leitura específica e pode
apresentar diferenças estruturais entre os arquivos, como presença ou
ausência de determinadas colunas. A função reduz o trabalho manual de
leitura, conferência e junção desses arquivos, automatizando a etapa
inicial de construção do banco.

\subsection{Uso}\label{uso}

\begin{Shaded}
\begin{Highlighting}[]
\NormalTok{dados\_RHC\_combinados }\OtherTok{\textless{}{-}} \FunctionTok{lerarquivoDBF}\NormalTok{()}
\end{Highlighting}
\end{Shaded}

\subsection{Entrada esperada}\label{entrada-esperada}

A função não exige argumentos externos. Para funcionar corretamente, os
arquivos DBF devem:

\begin{itemize}
\tightlist
\item
  Estar no diretório de trabalho (use \texttt{setwd()} para definir o
  diretório correto)
\item
  Manter o \textbf{nome original} proveniente do Integrador RHC, com o
  padrão iniciado por \texttt{rhc} (ex: \texttt{rhc20.dbf},
  \texttt{rhc21.dbf})
\end{itemize}

\subsection{Saída produzida}\label{sauxedda-produzida}

Um \texttt{data.frame} único contendo os registros combinados de todos
os arquivos DBF identificados. Caso nenhum arquivo válido seja
encontrado, a função retorna \texttt{NULL}.

\begin{center}\rule{0.5\linewidth}{0.5pt}\end{center}

\section{Comportamento operacional}\label{comportamento-operacional}

A função implementa um fluxo estruturado com os seguintes mecanismos:

\begin{description}
\tightlist
\item[\textbf{Verificação de dependências}]
O pacote \emph{foreign} é verificado e instalado automaticamente, se
necessário, garantindo execução autônoma.
\item[\textbf{Busca de arquivos}]
A função procura arquivos com o padrão
\texttt{\^{}rhc.*\textbackslash{}.dbf\$}. Se nenhum for encontrado,
encerra com mensagem informativa.
\item[\textbf{Processamento iterativo}]
Para cada arquivo, a função realiza leitura via
\texttt{foreign::read.dbf()}, verifica a presença de registros válidos,
e extrai o ano a partir do nome do arquivo.
\item[\textbf{Extração do ano}]
Utiliza expressão regular com lógica de conversão para século: anos
\texttt{\textless{}\ 50} são atribuídos ao século XXI; anos
\texttt{≥\ 50}, ao século XX.
\item[\textbf{Tratamento de erros}]
O uso de \texttt{tryCatch()} garante que erros individuais na leitura de
um arquivo não interrompam o processamento dos demais.
\item[\textbf{Harmonização estrutural}]
A função identifica o conjunto completo de colunas de todos os arquivos
e adiciona colunas ausentes com valores \texttt{NA}, compatibilizando
estruturas heterogêneas entre diferentes anos do RHC.
\item[\textbf{Integração final}]
Os dados são combinados via \texttt{do.call(rbind,\ ...)}, resultando em
um único \texttt{data.frame} consolidado.
\end{description}

\begin{tcolorbox}[enhanced jigsaw, arc=.35mm, bottomrule=.15mm, bottomtitle=1mm, breakable, colback=white, colbacktitle=quarto-callout-note-color!10!white, colframe=quarto-callout-note-color-frame, coltitle=black, left=2mm, leftrule=.75mm, opacityback=0, opacitybacktitle=0.6, rightrule=.15mm, title=\textcolor{quarto-callout-note-color}{\faInfo}\hspace{0.5em}{Importância no fluxo do RHCgen}, titlerule=0mm, toprule=.15mm, toptitle=1mm]

Esta é a \textbf{primeira etapa} do pipeline do RHCgen. Ela cria a base
inicial sobre a qual serão aplicadas as demais etapas de padronização,
recodificação, conversão de tipos, harmonização territorial, validação e
análise de qualidade dos dados.

\end{tcolorbox}

\bookmarksetup{startatroot}

\chapter{Padronização das
Variáveis}\label{padronizauxe7uxe3o-das-variuxe1veis}

\section{\texorpdfstring{Função
\texttt{renomear\_colunas()}}{Função renomear\_colunas()}}\label{funuxe7uxe3o-renomear_colunas}

\subsection{Finalidade}\label{finalidade-1}

A função \texttt{renomear\_colunas()} realiza a \textbf{padronização dos
nomes das variáveis} presentes no dataframe, substituindo as
nomenclaturas originais dos arquivos DBF por nomes mais descritivos,
consistentes e adequados ao uso analítico.

\subsection{Problema que resolve}\label{problema-que-resolve-1}

Os arquivos do RHC utilizam nomes de variáveis abreviados,
frequentemente pouco intuitivos e, em alguns casos, inconsistentes ao
longo do tempo. Essa característica dificulta a interpretação dos dados,
aumenta a probabilidade de erros na análise e compromete a
reprodutibilidade.

\subsection{Uso}\label{uso-1}

\begin{Shaded}
\begin{Highlighting}[]
\NormalTok{data }\OtherTok{\textless{}{-}} \FunctionTok{renomear\_colunas}\NormalTok{(data)}
\end{Highlighting}
\end{Shaded}

\subsection{Mapeamento de variáveis}\label{mapeamento-de-variuxe1veis}

O Tabela~\ref{tbl-mapeamento} apresenta o mapeamento completo das 47
variáveis originais do RHC para a nomenclatura padronizada adotada pelo
RHCgen.

\begin{longtable}[]{@{}ll@{}}
\caption{Mapeamento de variáveis do RHC para nomenclatura padronizada do
RHCgen (47 variáveis)}\label{tbl-mapeamento}\tabularnewline
\toprule\noalign{}
Nome original no DBF & Nome padronizado no RHCgen \\
\midrule\noalign{}
\endfirsthead
\toprule\noalign{}
Nome original no DBF & Nome padronizado no RHCgen \\
\midrule\noalign{}
\endhead
\bottomrule\noalign{}
\endlastfoot
\texttt{TPCASO} & \texttt{Tipo\_de\_Caso} \\
\texttt{SEXO} & \texttt{Sexo} \\
\texttt{IDADE} & \texttt{Idade} \\
\texttt{LOCALNAS} & \texttt{Local\_Nascimento} \\
\texttt{RACACOR} & \texttt{Raca\_Cor} \\
\texttt{INSTRUC} & \texttt{Escolaridade} \\
\texttt{CLIATEN} & \texttt{Clinica\_Atendimento} \\
\texttt{CLITRAT} & \texttt{Clinica\_Tratamento} \\
\texttt{HISTFAMC} & \texttt{Historico\_Familiar\_Cancer} \\
\texttt{ALCOOLIS} & \texttt{Consumo\_Alcool} \\
\texttt{TABAGISM} & \texttt{Tabagismo} \\
\texttt{ESTADRES} & \texttt{Estado\_Residencia} \\
\texttt{PROCEDEN} & \texttt{Codigo\_Municipio\_Residencia} \\
\texttt{ANOPRIDI} & \texttt{Ano\_Primeiro\_Diagnostico} \\
\texttt{ORIENC} & \texttt{Origem\_do\_Encaminhamento} \\
\texttt{EXDIAG} & \texttt{Exames\_Relevantes\_para\_Diagnostico} \\
\texttt{ESTCONJ} & \texttt{Estado\_conjugal} \\
\texttt{ANTRI} & \texttt{Ano\_Triagem} \\
\texttt{DTPRICON} & \texttt{Ano\_Primeira\_Consulta} \\
\texttt{DIAGANT} & \texttt{Diagnosticos\_e\_Tratamentos\_Anterior} \\
\texttt{BASMAIMP} & \texttt{Base\_Mais\_Importante\_Diagnostico} \\
\texttt{LOCTUDET} & \texttt{Localizacao\_Primaria\_3D} \\
\texttt{LOCTUPRI} & \texttt{Localizacao\_Primaria\_4D} \\
\texttt{TIPOHIST} & \texttt{Tipo\_Histologico} \\
\texttt{LATERALI} & \texttt{Lateralidade} \\
\texttt{LOCTUPRO} & \texttt{Local\_Provavel\_Tumor} \\
\texttt{MAISUMTU} & \texttt{Mais\_de\_Um\_Tumor} \\
\texttt{TNM} & \texttt{Classificacao\_TNM} \\
\texttt{ESTADIAM} & \texttt{Estadiamento\_Clinico} \\
\texttt{OUTROESTA} & \texttt{Outro\_Estadiamento\_Clinico} \\
\texttt{PTNM} & \texttt{TNM\_Patologico} \\
\texttt{RZNTR} & \texttt{Razao\_Nao\_Tratamento} \\
\texttt{DTINITRT} & \texttt{Ano\_Inicio\_Tratamento} \\
\texttt{PRITRATH} & \texttt{Primeiro\_Tratamento\_Hospital} \\
\texttt{ESTDFIMT} & \texttt{Estado\_Doenca\_Final\_Tratamento} \\
\texttt{DTTRIAGE} & \texttt{Data\_Triagem} \\
\texttt{DATAPRICON} & \texttt{Data\_Primeira\_Consulta} \\
\texttt{BASDIAGSP} &
\texttt{Base\_Mais\_Importante\_Diagnostico\_Sem\_Patologicas} \\
\texttt{CNES} & \texttt{CNES\_Hospital} \\
\texttt{UFUH} & \texttt{UF\_Unidade\_Hospital} \\
\texttt{MUUH} & \texttt{Municipio\_Unidade\_Hospital} \\
\texttt{OCUPACAO} & \texttt{Ocupacao\_Paciente} \\
\texttt{DTDIAGNO} & \texttt{Data\_Diagnostico} \\
\texttt{DATAINITRT} & \texttt{Data\_Inicio\_Primeiro\_Tratamento} \\
\texttt{DATAOBITO} & \texttt{Data\_Obito} \\
\texttt{VALOR\_TOT} & \texttt{Valor\_Total} \\
\texttt{ESTADIAG} &
\texttt{Estadiamento\_Clinico\_TNM\_grupo\_1985\_1999} \\
\end{longtable}

\begin{center}\rule{0.5\linewidth}{0.5pt}\end{center}

\section{\texorpdfstring{Função
\texttt{modificar\_tipo\_variavel()}}{Função modificar\_tipo\_variavel()}}\label{funuxe7uxe3o-modificar_tipo_variavel}

\subsection{Finalidade}\label{finalidade-2}

A função \texttt{modificar\_tipo\_variavel()} realiza a
\textbf{padronização dos tipos das variáveis}, convertendo
sistematicamente os campos para formatos apropriados ao processamento
analítico.

\subsection{Conversões realizadas}\label{conversuxf5es-realizadas}

{\def\LTcaptype{none} % do not increment counter
\begin{longtable}[]{@{}
  >{\raggedright\arraybackslash}p{(\linewidth - 2\tabcolsep) * \real{0.4865}}
  >{\raggedright\arraybackslash}p{(\linewidth - 2\tabcolsep) * \real{0.5135}}@{}}
\toprule\noalign{}
\begin{minipage}[b]{\linewidth}\raggedright
Tipo de variável
\end{minipage} & \begin{minipage}[b]{\linewidth}\raggedright
Formato resultante
\end{minipage} \\
\midrule\noalign{}
\endhead
\bottomrule\noalign{}
\endlastfoot
Variáveis categóricas em geral & \texttt{character} \\
\texttt{Idade} & \texttt{numeric} \\
Variáveis de data (\texttt{Data\_Triagem},
\texttt{Data\_Primeira\_Consulta}, \texttt{Data\_Diagnostico},
\texttt{Data\_Inicio\_Primeiro\_Tratamento}, \texttt{Data\_Obito}) &
\texttt{Date} (formato \texttt{\%d/\%m/\%Y}) \\
\end{longtable}
}

\subsection{Problema que resolve}\label{problema-que-resolve-2}

Arquivos oriundos do formato DBF frequentemente apresentam
inconsistências na tipagem após a importação para o R: datas lidas como
texto, variáveis numéricas tratadas como fatores, formatações
heterogêneas. Essas inconsistências podem gerar erros analíticos
silenciosos e comprometer a integridade das análises.

\subsection{Uso}\label{uso-2}

\begin{Shaded}
\begin{Highlighting}[]
\NormalTok{data }\OtherTok{\textless{}{-}} \FunctionTok{modificar\_tipo\_variavel}\NormalTok{(data)}
\end{Highlighting}
\end{Shaded}

\begin{tcolorbox}[enhanced jigsaw, arc=.35mm, bottomrule=.15mm, bottomtitle=1mm, breakable, colback=white, colbacktitle=quarto-callout-tip-color!10!white, colframe=quarto-callout-tip-color-frame, coltitle=black, left=2mm, leftrule=.75mm, opacityback=0, opacitybacktitle=0.6, rightrule=.15mm, title=\textcolor{quarto-callout-tip-color}{\faLightbulb}\hspace{0.5em}{Conversão segura de \texttt{Idade}}, titlerule=0mm, toprule=.15mm, toptitle=1mm]

A variável \texttt{Idade} é convertida em duas etapas: primeiro para
\texttt{character}, depois para \texttt{numeric}. Essa abordagem evita
problemas comuns associados à conversão direta de fatores para valores
numéricos, que pode introduzir distorções silenciosas nos dados.

\end{tcolorbox}

\bookmarksetup{startatroot}

\chapter{Recodificação das Variáveis
Categóricas}\label{recodificauxe7uxe3o-das-variuxe1veis-categuxf3ricas}

\section{\texorpdfstring{Função
\texttt{recodificar\_variaveis()}}{Função recodificar\_variaveis()}}\label{funuxe7uxe3o-recodificar_variaveis}

\subsection{Finalidade}\label{finalidade-3}

A função \texttt{recodificar\_variaveis()} realiza a
\textbf{transformação de variáveis categóricas} originalmente
codificadas em formato numérico ou textual para categorias
interpretáveis, com base no dicionário de variáveis do Registro
Hospitalar de Câncer. As variáveis são convertidas para o tipo
\texttt{factor} com níveis definidos, garantindo consistência semântica
e padronização para análise.

\subsection{Problema que resolve}\label{problema-que-resolve-3}

Os dados do RHC apresentam variáveis categóricas codificadas por números
(por exemplo, \texttt{"1"}, \texttt{"2"}, \texttt{"9"}), que não possuem
interpretação direta sem consulta ao dicionário. Essa codificação
dificulta a análise, aumenta o risco de erro na interpretação e
compromete a comunicação dos resultados.

\subsection{Uso}\label{uso-3}

\begin{Shaded}
\begin{Highlighting}[]
\NormalTok{data }\OtherTok{\textless{}{-}} \FunctionTok{recodificar\_variaveis}\NormalTok{(data)}
\end{Highlighting}
\end{Shaded}

\begin{tcolorbox}[enhanced jigsaw, arc=.35mm, bottomrule=.15mm, bottomtitle=1mm, breakable, colback=white, colbacktitle=quarto-callout-note-color!10!white, colframe=quarto-callout-note-color-frame, coltitle=black, left=2mm, leftrule=.75mm, opacityback=0, opacitybacktitle=0.6, rightrule=.15mm, title=\textcolor{quarto-callout-note-color}{\faInfo}\hspace{0.5em}{Tolerância a variações}, titlerule=0mm, toprule=.15mm, toptitle=1mm]

A função opera de forma tolerante a erros parciais: variáveis ausentes
são identificadas e reportadas sem interromper o fluxo. Valores não
previstos nas regras de recodificação são mantidos sem alteração.

\end{tcolorbox}

\begin{center}\rule{0.5\linewidth}{0.5pt}\end{center}

\section{Variáveis recodificadas}\label{variuxe1veis-recodificadas}

O Tabela~\ref{tbl-recodificacao} apresenta as variáveis categóricas
recodificadas, organizadas por domínio temático.

\begin{longtable}[]{@{}
  >{\raggedright\arraybackslash}p{(\linewidth - 4\tabcolsep) * \real{0.2093}}
  >{\raggedright\arraybackslash}p{(\linewidth - 4\tabcolsep) * \real{0.2326}}
  >{\raggedright\arraybackslash}p{(\linewidth - 4\tabcolsep) * \real{0.5581}}@{}}
\caption{Variáveis categóricas recodificadas pelo RHCgen, segundo
domínio temático}\label{tbl-recodificacao}\tabularnewline
\toprule\noalign{}
\begin{minipage}[b]{\linewidth}\raggedright
Domínio
\end{minipage} & \begin{minipage}[b]{\linewidth}\raggedright
Variável
\end{minipage} & \begin{minipage}[b]{\linewidth}\raggedright
Categorias padronizadas
\end{minipage} \\
\midrule\noalign{}
\endfirsthead
\toprule\noalign{}
\begin{minipage}[b]{\linewidth}\raggedright
Domínio
\end{minipage} & \begin{minipage}[b]{\linewidth}\raggedright
Variável
\end{minipage} & \begin{minipage}[b]{\linewidth}\raggedright
Categorias padronizadas
\end{minipage} \\
\midrule\noalign{}
\endhead
\bottomrule\noalign{}
\endlastfoot
\textbf{Identificação do caso} & \texttt{Tipo\_de\_Caso} & Analítico;
Não Analítico; Outro \\
\textbf{Sociodemográficas} & \texttt{Sexo} & Masculino; Feminino;
Outro \\
& \texttt{Raca\_Cor} & Branca; Preta; Amarela; Parda; Indígena; Sem
informação \\
& \texttt{Escolaridade} & Nenhuma; Fundamental incompleto; Fundamental
completo; Nível médio; Nível superior incompleto; Nível superior
completo; Sem informação \\
& \texttt{Estado\_conjugal} & Solteiro; Casado; Viúvo; Separado
judicialmente; União consensual; Sem informação \\
\textbf{Histórico e hábitos} & \texttt{Historico\_Familiar\_Cancer} &
Sim; Não; Sem informação \\
& \texttt{Consumo\_Alcool} & Nunca; Ex-consumidor; Sim; Não avaliado;
Não se aplica; Sem informação \\
& \texttt{Tabagismo} & Nunca; Ex-consumidor; Sim; Não avaliado; Não se
aplica; Sem informação \\
\textbf{Acesso e diagnóstico} & \texttt{Origem\_do\_Encaminhamento} &
SUS; Não SUS; Veio por conta própria; Não se aplica; Sem informação \\
& \texttt{Exames\_Relevantes\_para\_Diagnostico} & Exame clínico e
patologia clínica; Exames por imagem; Endoscopia e cirurgia exploradora;
Anatomia patológica; Marcadores tumorais; Não se aplica; Sem
informação \\
& \texttt{Base\_Mais\_Importante\_Diagnostico} & Clínica; Pesquisa
clínica; Exame por imagem; Marcadores tumorais; Citologia; Histologia da
metástase; Histologia do tumor primário; Sem informação \\
& \texttt{Base\_Mais\_Importante\_Diagnostico\_Sem\_Patologicas} & Exame
clínico; Recursos auxiliares não microscópicos; Confirmação
microscópica; Sem informação \\
\textbf{Características clínicas} & \texttt{Lateralidade} & Direita;
Esquerda; Bilateral; Não se aplica; Sem informação \\
& \texttt{Mais\_de\_Um\_Tumor} & Sim; Não; Duvidoso \\
& \texttt{Estadiamento\_Clinico\_TNM\_grupo\_1985\_1999} & Estádio I;
Estádio II; Estádio III; Estádio IV; Não se aplica; Sem informação \\
\textbf{Histórico assistencial} &
\texttt{Diagnosticos\_e\_Tratamentos\_Anterior} & Sem diagnóstico e sem
tratamento; Com diagnóstico e sem tratamento; Com diagnóstico e com
tratamento; Outros; Sem informação \\
\textbf{Tratamento} & \texttt{Primeiro\_Tratamento\_Hospital} & Nenhum;
Cirurgia; Radioterapia; Quimioterapia; Hormonioterapia; Transplante de
medula óssea; Imunoterapia; Outras; Sem informação \\
& \texttt{Razao\_Nao\_Tratamento} & Recusa do tratamento; Tratamento
realizado fora; Doença avançada ou comorbidades; Abandono; Complicações;
Óbito; Outras razões; Não se aplica; Sem informação \\
\textbf{Desfecho} & \texttt{Estado\_Doenca\_Final\_Tratamento} & Sem
evidência da doença; Remissão parcial; Doença estável; Doença em
progressão; Suporte terapêutico; Óbito; Não se aplica; Sem informação \\
\end{longtable}

Para consultar o dicionário completo de códigos, acesse:

\begin{quote}
\url{https://irhc.inca.gov.br/RHCNet/documentosTabWin.action?doc=dicionario.dados}
\end{quote}

\bookmarksetup{startatroot}

\chapter{Harmonização Territorial das Variáveis
Geográficas}\label{harmonizauxe7uxe3o-territorial-das-variuxe1veis-geogruxe1ficas}

\section{Funções de
harmonização}\label{funuxe7uxf5es-de-harmonizauxe7uxe3o}

O RHCgen dispõe de três funções dedicadas à conversão de códigos
geográficos em informações textuais interpretáveis:

{\def\LTcaptype{none} % do not increment counter
\begin{longtable}[]{@{}
  >{\raggedright\arraybackslash}p{(\linewidth - 2\tabcolsep) * \real{0.5714}}
  >{\raggedright\arraybackslash}p{(\linewidth - 2\tabcolsep) * \real{0.4286}}@{}}
\toprule\noalign{}
\begin{minipage}[b]{\linewidth}\raggedright
Função
\end{minipage} & \begin{minipage}[b]{\linewidth}\raggedright
Ação
\end{minipage} \\
\midrule\noalign{}
\endhead
\bottomrule\noalign{}
\endlastfoot
\texttt{renomear\_siglas\_estados()} & Converte siglas das UFs em nomes
completos \\
\texttt{renomear\_codigo\_municipio\_residencia()} & Converte código do
município de residência em nome \\
\texttt{renomear\_codigo\_municipio\_hospital()} & Converte código do
município do hospital em nome \\
\end{longtable}
}

\subsection{Uso}\label{uso-4}

\begin{Shaded}
\begin{Highlighting}[]
\NormalTok{data }\OtherTok{\textless{}{-}} \FunctionTok{renomear\_siglas\_estados}\NormalTok{(data)}
\NormalTok{data }\OtherTok{\textless{}{-}} \FunctionTok{renomear\_codigo\_municipio\_residencia}\NormalTok{(data)}
\NormalTok{data }\OtherTok{\textless{}{-}} \FunctionTok{renomear\_codigo\_municipio\_hospital}\NormalTok{(data)}
\end{Highlighting}
\end{Shaded}

\begin{center}\rule{0.5\linewidth}{0.5pt}\end{center}

\section{Finalidade das funções}\label{finalidade-das-funuxe7uxf5es}

As funções de harmonização territorial realizam a \textbf{conversão de
códigos geográficos} utilizados no RHC em informações textuais
interpretáveis, adicionando novas variáveis ao dataframe.
Especificamente, convertem siglas das Unidades da Federação e códigos de
municípios em seus respectivos nomes completos, tanto para o local de
residência do paciente quanto para a unidade hospitalar de atendimento.

\section{Problema que resolvem}\label{problema-que-resolvem}

Os dados do RHC utilizam codificações abreviadas para representação
geográfica --- siglas de estados e códigos numéricos de municípios.
Esses formatos dificultam a interpretação direta, limitam a análise
descritiva e comprometem a visualização espacial dos dados. As funções
resolvem esse problema ao incorporar, diretamente no banco, variáveis
com nomes completos e padronizados.

\section{Variáveis de entrada e
saída}\label{variuxe1veis-de-entrada-e-sauxedda}

{\def\LTcaptype{none} % do not increment counter
\begin{longtable}[]{@{}ll@{}}
\toprule\noalign{}
Variável de entrada & Variável criada \\
\midrule\noalign{}
\endhead
\bottomrule\noalign{}
\endlastfoot
\texttt{Estado\_Residencia} & \texttt{Nome\_Estado\_Residencia} \\
\texttt{UF\_Unidade\_Hospital} & \texttt{Nome\_Estado\_Hospital} \\
\texttt{Codigo\_Municipio\_Residencia} &
\texttt{Nome\_Municipio\_Residencia} \\
\texttt{Municipio\_Unidade\_Hospital} &
\texttt{Nome\_Municipio\_Hospital} \\
\end{longtable}
}

\begin{tcolorbox}[enhanced jigsaw, arc=.35mm, bottomrule=.15mm, bottomtitle=1mm, breakable, colback=white, colbacktitle=quarto-callout-note-color!10!white, colframe=quarto-callout-note-color-frame, coltitle=black, left=2mm, leftrule=.75mm, opacityback=0, opacitybacktitle=0.6, rightrule=.15mm, title=\textcolor{quarto-callout-note-color}{\faInfo}\hspace{0.5em}{Preservação das variáveis originais}, titlerule=0mm, toprule=.15mm, toptitle=1mm]

As funções de harmonização territorial \textbf{preservam os códigos
originais}, adicionando novas colunas ao dataframe sem alterá-las.
Valores não mapeados resultam em \texttt{NA}, sem interrupção do fluxo
de execução.

\end{tcolorbox}

\section{Características
operacionais}\label{caracteruxedsticas-operacionais}

\begin{itemize}
\tightlist
\item
  Operam de forma \textbf{vetorizada}, garantindo eficiência em grandes
  volumes de dados
\item
  Implementam \textbf{verificação explícita} da existência das variáveis
  de entrada
\item
  Utilizam interrupção controlada (\texttt{stop()}) quando condições
  mínimas não são atendidas
\item
  Não validam a consistência dos códigos de entrada, assumindo
  conformidade com o padrão do RHC
\end{itemize}

\begin{tcolorbox}[enhanced jigsaw, arc=.35mm, bottomrule=.15mm, bottomtitle=1mm, breakable, colback=white, colbacktitle=quarto-callout-tip-color!10!white, colframe=quarto-callout-tip-color-frame, coltitle=black, left=2mm, leftrule=.75mm, opacityback=0, opacitybacktitle=0.6, rightrule=.15mm, title=\textcolor{quarto-callout-tip-color}{\faLightbulb}\hspace{0.5em}{Aplicações analíticas}, titlerule=0mm, toprule=.15mm, toptitle=1mm]

Essa etapa é fundamental para análises espaciais, estudos de
acessibilidade geográfica, avaliação de fluxos de deslocamento de
pacientes e construção de indicadores territoriais. A presença de nomes
completos facilita a geração de mapas e a comunicação dos resultados em
contextos científicos e institucionais.

\end{tcolorbox}

\bookmarksetup{startatroot}

\chapter{Padronização de Códigos Clínicos e
Assistenciais}\label{padronizauxe7uxe3o-de-cuxf3digos-cluxednicos-e-assistenciais}

\section{Códigos diagnósticos}\label{cuxf3digos-diagnuxf3sticos}

\subsection{Funções disponíveis}\label{funuxe7uxf5es-disponuxedveis}

{\def\LTcaptype{none} % do not increment counter
\begin{longtable}[]{@{}
  >{\raggedright\arraybackslash}p{(\linewidth - 2\tabcolsep) * \real{0.5714}}
  >{\raggedright\arraybackslash}p{(\linewidth - 2\tabcolsep) * \real{0.4286}}@{}}
\toprule\noalign{}
\begin{minipage}[b]{\linewidth}\raggedright
Função
\end{minipage} & \begin{minipage}[b]{\linewidth}\raggedright
Ação
\end{minipage} \\
\midrule\noalign{}
\endhead
\bottomrule\noalign{}
\endlastfoot
\texttt{renomear\_CID\_3digitos()} & CID-O com 3 dígitos → descrição
completa da localização tumoral \\
\texttt{renomear\_CID\_4digitos()} & CID-O com 4 dígitos → descrição
completa da localização tumoral \\
\texttt{renomear\_tipo\_histologico()} & Código histológico → descrição
do tipo histológico \\
\texttt{renomear\_estadiamento\_clinico()} & Código de estadiamento →
categoria padronizada \\
\texttt{classificacao\_CID\_0()} & CID-O → agrupamento por grande grupo
anatômico \\
\end{longtable}
}

\subsection{Uso}\label{uso-5}

\begin{Shaded}
\begin{Highlighting}[]
\NormalTok{data }\OtherTok{\textless{}{-}} \FunctionTok{renomear\_CID\_3digitos}\NormalTok{(data)}
\NormalTok{data }\OtherTok{\textless{}{-}} \FunctionTok{renomear\_tipo\_histologico}\NormalTok{(data)}
\NormalTok{data }\OtherTok{\textless{}{-}} \FunctionTok{renomear\_estadiamento\_clinico}\NormalTok{(data)}
\NormalTok{data }\OtherTok{\textless{}{-}} \FunctionTok{classificacao\_CID\_0}\NormalTok{(data)}
\end{Highlighting}
\end{Shaded}

\subsection{Variáveis criadas}\label{variuxe1veis-criadas}

{\def\LTcaptype{none} % do not increment counter
\begin{longtable}[]{@{}ll@{}}
\toprule\noalign{}
Variável de entrada & Variável criada \\
\midrule\noalign{}
\endhead
\bottomrule\noalign{}
\endlastfoot
\texttt{Localizacao\_Primaria\_3D} & \texttt{CID3d} \\
\texttt{Localizacao\_Primaria\_4D} & \texttt{CID4d} \\
\texttt{Tipo\_Histologico} & \texttt{Tipo\_Histologico\_Completo} \\
\texttt{Estadiamento\_Clinico} & \texttt{Nome\_Estadiamento\_Clinico} \\
\texttt{Localizacao\_Primaria\_3D} & \texttt{Classificacao\_CID\_O} \\
\end{longtable}
}

\begin{center}\rule{0.5\linewidth}{0.5pt}\end{center}

\section{\texorpdfstring{Classificação por grupo anatômico:
\texttt{classificacao\_CID\_0()}}{Classificação por grupo anatômico: classificacao\_CID\_0()}}\label{classificauxe7uxe3o-por-grupo-anatuxf4mico-classificacao_cid_0}

\subsection{Finalidade}\label{finalidade-4}

A função \texttt{classificacao\_CID\_0()} classifica as neoplasias
malignas em \textbf{grandes grupos anatômicos} a partir dos códigos
registrados na variável \texttt{Localizacao\_Primaria\_3D}. A variável
resultante, \texttt{Classificacao\_CID\_O}, agrupa os códigos CID-O em
categorias como neoplasias malignas dos órgãos digestivos, órgãos
respiratórios, mama, órgãos genitais masculinos, órgãos urinários e
tecidos linfático/hematopoético.

\subsection{Lógica de
classificação}\label{luxf3gica-de-classificauxe7uxe3o}

Para cada registro, a função aplica três decisões sequenciais:

\begin{enumerate}
\def\labelenumi{\arabic{enumi}.}
\tightlist
\item
  Se o código é \texttt{NA} → resultado permanece \texttt{NA}
\item
  Se o código está no vetor de mapeamento → recebe a categoria
  correspondente
\item
  Se o código não está previsto → recebe \texttt{"Outros"}
\end{enumerate}

\subsection{Padronização do estadiamento
clínico}\label{padronizauxe7uxe3o-do-estadiamento-cluxednico}

A função \texttt{renomear\_estadiamento\_clinico()} resolve um problema
crítico de heterogeneidade: o RHC acumula múltiplos formatos de registro
do estadiamento ao longo do tempo (por exemplo: \texttt{"1"},
\texttt{"I"}, \texttt{"1A"}, \texttt{"01"}). A função mapeia todos esses
formatos para categorias padronizadas (\texttt{"Estádio\ I"},
\texttt{"Estádio\ II"}, etc.), promovendo comparabilidade entre
registros de diferentes períodos.

\begin{center}\rule{0.5\linewidth}{0.5pt}\end{center}

\section{Códigos assistenciais e
hospitalares}\label{cuxf3digos-assistenciais-e-hospitalares}

\subsection{Funções disponíveis}\label{funuxe7uxf5es-disponuxedveis-1}

{\def\LTcaptype{none} % do not increment counter
\begin{longtable}[]{@{}
  >{\raggedright\arraybackslash}p{(\linewidth - 2\tabcolsep) * \real{0.5714}}
  >{\raggedright\arraybackslash}p{(\linewidth - 2\tabcolsep) * \real{0.4286}}@{}}
\toprule\noalign{}
\begin{minipage}[b]{\linewidth}\raggedright
Função
\end{minipage} & \begin{minipage}[b]{\linewidth}\raggedright
Ação
\end{minipage} \\
\midrule\noalign{}
\endhead
\bottomrule\noalign{}
\endlastfoot
\texttt{renomear\_clinica()} & Converte códigos de clínicas para nomes
completos \\
\texttt{renomear\_CNES()} & Converte código CNES para nome do
estabelecimento \\
\end{longtable}
}

\subsection{Uso}\label{uso-6}

\begin{Shaded}
\begin{Highlighting}[]
\NormalTok{data }\OtherTok{\textless{}{-}} \FunctionTok{renomear\_clinica}\NormalTok{(data)}
\NormalTok{data }\OtherTok{\textless{}{-}} \FunctionTok{renomear\_CNES}\NormalTok{(data)}
\end{Highlighting}
\end{Shaded}

\subsection{Variáveis criadas}\label{variuxe1veis-criadas-1}

{\def\LTcaptype{none} % do not increment counter
\begin{longtable}[]{@{}ll@{}}
\toprule\noalign{}
Variável de entrada & Variável criada \\
\midrule\noalign{}
\endhead
\bottomrule\noalign{}
\endlastfoot
\texttt{Clinica\_Atendimento} & \texttt{Nome\_Clinica\_Atendimento} \\
\texttt{Clinica\_Tratamento} & \texttt{Nome\_Clinica\_Tratamento} \\
\texttt{CNES\_Hospital} & \texttt{Estabelecimento\_Hospitalar} \\
\end{longtable}
}

Para consultar o dicionário de códigos de clínicas:

\begin{quote}
\url{https://irhc.inca.gov.br/RHCNet/documentosTabWin.action?doc=tabela.clinicas}
\end{quote}

\begin{center}\rule{0.5\linewidth}{0.5pt}\end{center}

\section{\texorpdfstring{Organização estrutural do banco:
\texttt{ordenando\_colunas()}}{Organização estrutural do banco: ordenando\_colunas()}}\label{organizauxe7uxe3o-estrutural-do-banco-ordenando_colunas}

Após todas as etapas de leitura, padronização, recodificação e criação
de novas variáveis, a função \texttt{ordenando\_colunas()} organiza as
variáveis do dataframe em uma \textbf{ordem lógica predefinida},
alinhada à estrutura analítica do RHCgen.

\begin{Shaded}
\begin{Highlighting}[]
\NormalTok{data }\OtherTok{\textless{}{-}} \FunctionTok{ordenando\_colunas}\NormalTok{(data)}
\end{Highlighting}
\end{Shaded}

A organização segue uma sequência que integra dimensões
sociodemográficas, clínicas, assistenciais e temporais, facilitando a
utilização do banco em estudos epidemiológicos.

\begin{tcolorbox}[enhanced jigsaw, arc=.35mm, bottomrule=.15mm, bottomtitle=1mm, breakable, colback=white, colbacktitle=quarto-callout-note-color!10!white, colframe=quarto-callout-note-color-frame, coltitle=black, left=2mm, leftrule=.75mm, opacityback=0, opacitybacktitle=0.6, rightrule=.15mm, title=\textcolor{quarto-callout-note-color}{\faInfo}\hspace{0.5em}{Tolerância a variações estruturais}, titlerule=0mm, toprule=.15mm, toptitle=1mm]

A função não adiciona colunas ausentes nem interrompe a execução quando
variáveis não estão disponíveis, operando de forma tolerante à variação
estrutural do banco entre diferentes períodos históricos do RHC.

\end{tcolorbox}

\bookmarksetup{startatroot}

\chapter{Avaliação da Qualidade dos
Dados}\label{avaliauxe7uxe3o-da-qualidade-dos-dados}

\section{\texorpdfstring{Função principal:
\texttt{analise\_qualidade\_dos\_dados()}}{Função principal: analise\_qualidade\_dos\_dados()}}\label{funuxe7uxe3o-principal-analise_qualidade_dos_dados}

\subsection{Finalidade}\label{finalidade-5}

A função \texttt{analise\_qualidade\_dos\_dados()} executa uma
\textbf{avaliação automatizada da qualidade} do banco processado pelo
RHCgen. Essa etapa reúne análises de completude e consistência,
permitindo identificar variáveis com maior proporção de dados ausentes,
códigos inválidos, idades inconsistentes, inconsistências temporais e
incompatibilidades entre sexo e topografia tumoral.

\subsection{Uso}\label{uso-7}

\begin{Shaded}
\begin{Highlighting}[]
\FunctionTok{analise\_qualidade\_dos\_dados}\NormalTok{(data)}
\end{Highlighting}
\end{Shaded}

\begin{tcolorbox}[enhanced jigsaw, arc=.35mm, bottomrule=.15mm, bottomtitle=1mm, breakable, colback=white, colbacktitle=quarto-callout-note-color!10!white, colframe=quarto-callout-note-color-frame, coltitle=black, left=2mm, leftrule=.75mm, opacityback=0, opacitybacktitle=0.6, rightrule=.15mm, title=\textcolor{quarto-callout-note-color}{\faInfo}\hspace{0.5em}{Saídas produzidas}, titlerule=0mm, toprule=.15mm, toptitle=1mm]

A função não retorna um novo dataframe. Ela exibe mensagens de progresso
no console, gera um gráfico de barras de completude e produz um
\textbf{relatório HTML} com as verificações de inconsistência, aberto
automaticamente no navegador ao final da execução.

\end{tcolorbox}

\begin{center}\rule{0.5\linewidth}{0.5pt}\end{center}

\section{Análise de completude}\label{anuxe1lise-de-completude}

\subsection{\texorpdfstring{Função
\texttt{analise\_completude()}}{Função analise\_completude()}}\label{funuxe7uxe3o-analise_completude}

Para cada variável do banco, a função calcula o número de valores
ausentes, considerando:

\begin{itemize}
\tightlist
\item
  Valores \texttt{NA}
\item
  Registros explicitamente classificados como \texttt{"Sem\ informação"}
\end{itemize}

\subsubsection{Classificação da
completude}\label{classificauxe7uxe3o-da-completude}

A classificação segue o escore proposto por Romero e Cunha (2006):

{\def\LTcaptype{none} % do not increment counter
\begin{longtable}[]{@{}ll@{}}
\toprule\noalign{}
Percentual de completude & Classificação \\
\midrule\noalign{}
\endhead
\bottomrule\noalign{}
\endlastfoot
≥ 95\% & Excelente \\
90\% a \textless{} 95\% & Bom \\
80\% a \textless{} 90\% & Regular \\
50\% a \textless{} 80\% & Ruim \\
\textless{} 50\% & Muito Ruim \\
\end{longtable}
}

\begin{tcolorbox}[enhanced jigsaw, arc=.35mm, bottomrule=.15mm, bottomtitle=1mm, breakable, colback=white, colbacktitle=quarto-callout-tip-color!10!white, colframe=quarto-callout-tip-color-frame, coltitle=black, left=2mm, leftrule=.75mm, opacityback=0, opacitybacktitle=0.6, rightrule=.15mm, title=\textcolor{quarto-callout-tip-color}{\faLightbulb}\hspace{0.5em}{Ajuste para \texttt{Data\_Obito}}, titlerule=0mm, toprule=.15mm, toptitle=1mm]

Quando a variável \texttt{Data\_Obito} está presente, a completude é
ajustada considerando apenas registros em que o desfecho foi óbito,
permitindo uma avaliação mais adequada dessa variável no contexto
clínico.

\end{tcolorbox}

\subsection{\texorpdfstring{Análise de completude por ano:
\texttt{analise\_completude\_ano()}}{Análise de completude por ano: analise\_completude\_ano()}}\label{anuxe1lise-de-completude-por-ano-analise_completude_ano}

A função \texttt{analise\_completude\_ano()} calcula a completude de
cada variável estratificada por ano do banco:

\begin{Shaded}
\begin{Highlighting}[]
\NormalTok{tabela\_completude\_ano }\OtherTok{\textless{}{-}} \FunctionTok{analise\_completude\_ano}\NormalTok{(dados\_RHC\_combinados)}
\end{Highlighting}
\end{Shaded}

O resultado é um dataframe no qual cada linha representa uma variável e
cada coluna corresponde a um ano. As células apresentam o percentual e a
classificação abreviada, por exemplo: \texttt{95.2\%\ (E)}.

{\def\LTcaptype{none} % do not increment counter
\begin{longtable}[]{@{}ll@{}}
\toprule\noalign{}
Sigla & Classificação \\
\midrule\noalign{}
\endhead
\bottomrule\noalign{}
\endlastfoot
E & Excelente (≥ 95\%) \\
B & Bom (90\% a \textless{} 95\%) \\
R & Regular (80\% a \textless{} 90\%) \\
RU & Ruim (50\% a \textless{} 80\%) \\
MR & Muito Ruim (\textless{} 50\%) \\
\end{longtable}
}

\begin{center}\rule{0.5\linewidth}{0.5pt}\end{center}

\section{Análise de consistência}\label{anuxe1lise-de-consistuxeancia}

\subsection{\texorpdfstring{Função
\texttt{tabela\_inconsistencias()}}{Função tabela\_inconsistencias()}}\label{funuxe7uxe3o-tabela_inconsistencias}

A análise de consistência é organizada em duas dimensões:

\subsubsection{1. Validação de códigos e preenchimento
formal}\label{validauxe7uxe3o-de-cuxf3digos-e-preenchimento-formal}

\textbf{Função: \texttt{verificar\_codigos\_inconsistentes()}}

Avalia se os valores registrados em variáveis selecionadas pertencem ao
conjunto de códigos ou categorias esperadas para cada campo, incluindo:
\texttt{Sexo}, \texttt{Raca\_Cor}, \texttt{Escolaridade},
\texttt{Localizacao\_Primaria\_3D}, \texttt{Localizacao\_Primaria\_4D},
\texttt{Tipo\_Histologico}, \texttt{CNES\_Hospital} e
\texttt{UF\_Unidade\_Hospital}.

\begin{tcolorbox}[enhanced jigsaw, arc=.35mm, bottomrule=.15mm, bottomtitle=1mm, breakable, colback=white, colbacktitle=quarto-callout-note-color!10!white, colframe=quarto-callout-note-color-frame, coltitle=black, left=2mm, leftrule=.75mm, opacityback=0, opacitybacktitle=0.6, rightrule=.15mm, title=\textcolor{quarto-callout-note-color}{\faInfo}\hspace{0.5em}{Nota}, titlerule=0mm, toprule=.15mm, toptitle=1mm]

Códigos como \texttt{99}, \texttt{999}, \texttt{88} e \texttt{888},
quando previstos no dicionário interno, são tratados como
\textbf{categorias válidas}, e não como inconsistências.

\end{tcolorbox}

\textbf{Função: \texttt{verificar\_idades\_inconsistentes()}}

Avalia a plausibilidade da variável \texttt{Idade}, aplicando dois
critérios: idade menor que 0 e idade maior que 150 anos.

\subsubsection{2. Coerência lógico-clínica entre
variáveis}\label{coeruxeancia-luxf3gico-cluxednica-entre-variuxe1veis}

\textbf{Função: \texttt{verificar\_anos\_inconsistentes()}}

Avalia a coerência cronológica entre anos dos principais eventos
clínicos e assistenciais, verificando cinco comparações lógicas:

\begin{enumerate}
\def\labelenumi{\arabic{enumi}.}
\tightlist
\item
  Ano do diagnóstico anterior à primeira consulta
\item
  Ano do diagnóstico anterior à triagem
\item
  Início do tratamento anterior à primeira consulta
\item
  Início do tratamento anterior à triagem
\item
  Início do tratamento anterior ao primeiro diagnóstico
\end{enumerate}

\textbf{Função: \texttt{verificar\_cancer\_sexo\_inconsistentes()}}

Avalia a compatibilidade entre \texttt{Sexo} e a localização primária do
tumor para topografias anatomicamente relacionadas ao sexo (órgãos
genitais femininos e masculinos).

\begin{tcolorbox}[enhanced jigsaw, arc=.35mm, bottomrule=.15mm, bottomtitle=1mm, breakable, colback=white, colbacktitle=quarto-callout-important-color!10!white, colframe=quarto-callout-important-color-frame, coltitle=black, left=2mm, leftrule=.75mm, opacityback=0, opacitybacktitle=0.6, rightrule=.15mm, title=\textcolor{quarto-callout-important-color}{\faExclamation}\hspace{0.5em}{Denominador específico}, titlerule=0mm, toprule=.15mm, toptitle=1mm]

A proporção de inconsistência sexo--topografia é calculada usando como
denominador apenas os registros completos pertencentes ao respectivo
grupo topográfico, e não o total do banco, evitando distorção da
estimativa.

\end{tcolorbox}

\subsection{Classificação de
gravidade}\label{classificauxe7uxe3o-de-gravidade}

{\def\LTcaptype{none} % do not increment counter
\begin{longtable}[]{@{}ll@{}}
\toprule\noalign{}
Proporção de inconsistência & Classificação \\
\midrule\noalign{}
\endhead
\bottomrule\noalign{}
\endlastfoot
\textless{} 1\% & Baixo \\
1\% a \textless{} 5\% & Moderado \\
5\% a \textless{} 20\% & Alto \\
≥ 20\% & Crítico \\
\end{longtable}
}

\subsection{Relatório HTML}\label{relatuxf3rio-html}

Os resultados são integrados em um relatório HTML estruturado,
incluindo:

\begin{itemize}
\tightlist
\item
  Resumo executivo com principais inconsistências
\item
  Tabelas ordenadas por magnitude do problema
\item
  Classificação de gravidade com codificação visual
\item
  Interpretação automatizada por seção
\item
  Notas metodológicas específicas para cada análise
\item
  Seção de implicações analíticas e limitações
\end{itemize}

O relatório é salvo automaticamente no diretório de trabalho e aberto no
navegador padrão ao final da execução.

\bookmarksetup{startatroot}

\chapter{Aplicação}\label{aplicauxe7uxe3o}

\section{Processamento da base nacional
(2000--2023)}\label{processamento-da-base-nacional-20002023}

Para demonstrar o funcionamento do RHCgen em escala real, foram
processados \textbf{24 arquivos do RHC}, correspondentes ao período de
2000 a 2023, disponibilizados pelo IntegradorRHC.

O processamento foi realizado por meio da função
\texttt{construir\_banco()}, que integrou automaticamente os dados
provenientes de diferentes anos e promoveu a padronização estrutural e
semântica das variáveis.

\section{Resultados operacionais}\label{resultados-operacionais}

{\def\LTcaptype{none} % do not increment counter
\begin{longtable}[]{@{}ll@{}}
\toprule\noalign{}
Indicador & Valor \\
\midrule\noalign{}
\endhead
\bottomrule\noalign{}
\endlastfoot
Arquivos processados & 24 (anos 2000--2023) \\
Registros na base final & \textbf{5.394.771} \\
Variáveis na base final & \textbf{59} \\
Tempo de execução & \textbf{4,61 minutos} \\
Consumo de memória RAM & \textbf{≈ 2,01 GB} \\
\end{longtable}
}

\section{Código de reprodução}\label{cuxf3digo-de-reproduuxe7uxe3o}

\begin{Shaded}
\begin{Highlighting}[]
\FunctionTok{library}\NormalTok{(RHCgen)}

\CommentTok{\# Definir o diretório com os arquivos .dbf}
\FunctionTok{setwd}\NormalTok{(}\StringTok{"caminho/para/seus/arquivos/dbf"}\NormalTok{)}

\CommentTok{\# Executar o pipeline completo}
\NormalTok{dados\_RHC }\OtherTok{\textless{}{-}} \FunctionTok{construir\_banco}\NormalTok{()}

\CommentTok{\# Verificar estrutura da base}
\FunctionTok{dim}\NormalTok{(dados\_RHC)       }\CommentTok{\# 5394771 linhas × 59 colunas}
\FunctionTok{str}\NormalTok{(dados\_RHC)       }\CommentTok{\# estrutura das variáveis}
\FunctionTok{head}\NormalTok{(dados\_RHC)      }\CommentTok{\# primeiros registros}

\CommentTok{\# Análise de completude por ano (opcional)}
\NormalTok{tabela\_completude }\OtherTok{\textless{}{-}} \FunctionTok{analise\_completude\_ano}\NormalTok{(dados\_RHC)}
\end{Highlighting}
\end{Shaded}

\begin{tcolorbox}[enhanced jigsaw, arc=.35mm, bottomrule=.15mm, bottomtitle=1mm, breakable, colback=white, colbacktitle=quarto-callout-tip-color!10!white, colframe=quarto-callout-tip-color-frame, coltitle=black, left=2mm, leftrule=.75mm, opacityback=0, opacitybacktitle=0.6, rightrule=.15mm, title=\textcolor{quarto-callout-tip-color}{\faLightbulb}\hspace{0.5em}{Exportação da base}, titlerule=0mm, toprule=.15mm, toptitle=1mm]

O RHCgen não salva automaticamente os dados em disco. Para exportar a
base analítica, utilize, por exemplo:

\begin{Shaded}
\begin{Highlighting}[]
\CommentTok{\# Formato R nativo (recomendado para reutilização)}
\FunctionTok{saveRDS}\NormalTok{(dados\_RHC, }\StringTok{"dados\_RHC\_2000\_2023.rds"}\NormalTok{)}

\CommentTok{\# Formato CSV (maior compatibilidade)}
\FunctionTok{write.csv}\NormalTok{(dados\_RHC, }\StringTok{"dados\_RHC\_2000\_2023.csv"}\NormalTok{, }\AttributeTok{row.names =} \ConstantTok{FALSE}\NormalTok{)}
\end{Highlighting}
\end{Shaded}

\end{tcolorbox}

\section{Considerações sobre requisitos
computacionais}\label{considerauxe7uxf5es-sobre-requisitos-computacionais}

O consumo de \textbf{2,01 GB de memória RAM} para a base completa
(2000--2023) pode representar uma limitação em ambientes computacionais
com recursos reduzidos. Para usuários que trabalhem com infraestrutura
restrita, recomenda-se:

\begin{itemize}
\tightlist
\item
  Processar subconjuntos de arquivos por período (ex: décadas separadas)
\item
  Aumentar a memória disponível para o R com \texttt{memory.limit()}
  (Windows) ou configurações de sistema
\item
  Considerar o uso de pacotes como \emph{data.table} ou \emph{arrow}
  para manipulação de grandes volumes após o processamento inicial
\end{itemize}

\bookmarksetup{startatroot}

\chapter{Discussão}\label{discussuxe3o}

O RHCgen demonstrou ser uma ferramenta eficiente para a automação e
padronização do processamento de dados do Registro Hospitalar de Câncer
no Brasil, ao consolidar 5.394.771 registros provenientes de 24 arquivos
anuais (2000--2023) em uma única base analítica estruturada em 4,61
minutos e com consumo de aproximadamente 2,01 GB de memória. Esse
resultado evidencia a viabilidade técnica do pacote para o manejo de
grandes volumes de dados epidemiológicos em saúde, corroborando o
potencial das ferramentas de ciência de dados de código aberto para
transformar fluxos de trabalho em pesquisa em saúde pública ({Lowndes et
al.} 2017; Chan 2018).

\section{Vantagens em relação ao
TabWin}\label{vantagens-em-relauxe7uxe3o-ao-tabwin}

A utilização da linguagem R como plataforma de desenvolvimento
mostrou-se adequada para os objetivos propostos. Diferentemente do
\textbf{TabWin} --- ferramenta amplamente utilizada no contexto do
DATASUS, mas restrita ao sistema operacional Windows e com capacidades
analíticas limitadas ---, o RHCgen:

\begin{itemize}
\tightlist
\item
  Opera em \textbf{ambiente multiplataforma} (Windows, Linux, macOS)
\item
  Integra-se nativamente ao \textbf{ecossistema estatístico do R}
\item
  Elimina a necessidade de exportação dos dados para análises
  complementares
\end{itemize}

Esse aspecto é relevante porque a dependência de etapas manuais de
exportação e recodificação constitui uma fonte recorrente de erros e
inconsistências no processamento de sistemas de informação em saúde no
Brasil (Jorge et al. 2010).

\section{Posicionamento no ecossistema de pacotes para dados do
SUS}\label{posicionamento-no-ecossistema-de-pacotes-para-dados-do-sus}

Nesse sentido, o RHCgen se alinha a iniciativas brasileiras anteriores
voltadas à superação dessas limitações, como o \emph{microdatasus}
(Saldanha et al. 2019) e o \emph{Rtabnetsp} (Morais et al. 2021),
ampliando o conjunto de ferramentas disponíveis para o processamento de
dados do SUS em R.

\section{Contribuição para a qualidade dos
dados}\label{contribuiuxe7uxe3o-para-a-qualidade-dos-dados}

O RHCgen, ao automatizar a análise de completude e classificar as
variáveis segundo critérios validados, oferece ao pesquisador um
\textbf{diagnóstico imediato da qualidade do banco}, favorecendo
decisões informadas sobre:

\begin{itemize}
\tightlist
\item
  A viabilidade de determinadas análises
\item
  A necessidade de estratégias de imputação ou restrição de subgrupos
\item
  A comparabilidade entre períodos históricos
\end{itemize}

Mudanças estruturais nas fichas do RHC ao longo do período analisado
podem contribuir para a incompatibilidade entre registros históricos e
regras de validação mais recentes. Esse padrão reflete a heterogeneidade
estrutural dos arquivos anuais --- um dos principais desafios
identificados no desenvolvimento do RHCgen --- e reforça a necessidade
de ferramentas que realizem harmonização das bases, como a função
\texttt{lerarquivoDBF()}.

O fato de o RHCgen \textbf{quantificar e estratificar inconsistências
temporais por tipo de verificação} representa um avanço metodológico
concreto em relação ao processamento convencional, que em geral não
contempla essa camada de validação de forma automatizada e reprodutível.

\section{Ciência aberta e
reprodutibilidade}\label{ciuxeancia-aberta-e-reprodutibilidade}

Do ponto de vista da contribuição para a ciência aberta e para a
reprodutibilidade da pesquisa em oncologia, o RHCgen alinha-se a um
movimento mais amplo de valorização de ferramentas de código aberto no
contexto da saúde pública ({Lowndes et al.} 2017; Jalal et al. 2017). A
disponibilização do pacote no GitHub, acompanhada de guia de uso
acessível, \textbf{reduz a barreira de entrada para pesquisadores} sem
formação específica em ciência de dados, democratizando o acesso a uma
base de dados com mais de cinco milhões de registros oncológicos de
relevância nacional.

Iniciativas análogas, como o \emph{microdatasus} e o \emph{Rtabnetsp},
demonstraram que pacotes R voltados a dados do SUS têm ampla adoção pela
comunidade científica brasileira (Saldanha et al. 2019; Morais et al.
2021), o que sugere potencial similar para o RHCgen.

\section{Limitações}\label{limitauxe7uxf5es}

Algumas limitações do estudo merecem ser reconhecidas:

\begin{description}
\tightlist
\item[\textbf{Download manual dos arquivos}]
A plataforma IRHC não disponibiliza acesso remoto automatizado via FTP
ou API, impondo uma etapa não automatizável que pode constituir
obstáculo para usuários menos familiarizados com o processo de coleta de
dados.
\item[\textbf{Dados recentes em fase de validação}]
Os registros mais recentes do IRHC podem ainda estar em fase de
validação, com maior proporção de campos incompletos ou revisados, o que
pode afetar a qualidade de análises baseadas nos anos mais próximos da
data de extração.
\item[\textbf{Requisitos de memória}]
O consumo de 2,01 GB de memória RAM pode representar limitação em
ambientes computacionais com recursos reduzidos.
\end{description}

\bookmarksetup{startatroot}

\chapter{Conclusão}\label{conclusuxe3o}

O RHCgen constitui uma contribuição metodológica para o campo da
epidemiologia oncológica no Brasil, ao oferecer uma solução
\textbf{padronizada, reprodutível e de acesso aberto} para o
processamento de dados do Registro Hospitalar de Câncer.

A capacidade de consolidar mais de cinco milhões de registros
distribuídos ao longo de 24 anos em uma única base analítica
estruturada, por meio de um fluxo automatizado e auditável, representa
um avanço concreto em relação aos métodos convencionais de manipulação
manual de dados, historicamente sujeitos a inconsistências e de difícil
reprodução entre diferentes equipes de pesquisa.

\section{Contribuições técnicas}\label{contribuiuxe7uxf5es-tuxe9cnicas}

Do ponto de vista técnico, o pacote demonstrou:

\begin{itemize}
\tightlist
\item
  \textbf{Desempenho satisfatório} em termos de tempo de execução e
  integridade das transformações aplicadas, validadas contra os
  dicionários oficiais do SisRHC
\item
  \textbf{Funções modulares} que permitem ao usuário tanto a execução
  integrada do fluxo completo quanto o controle granular de etapas
  específicas, adaptando-se a diferentes demandas analíticas e níveis de
  experiência com dados do SUS
\item
  \textbf{Diagnóstico imediato de qualidade}, com análises de completude
  e consistência automaticamente geradas --- elemento indispensável para
  decisões metodológicas sobre a viabilidade de análises e a adoção de
  estratégias de tratamento de dados faltantes
\end{itemize}

\section{Impacto esperado}\label{impacto-esperado}

A disponibilização pública do RHCgen no GitHub, acompanhada de
documentação acessível, posiciona a ferramenta como um \textbf{recurso
colaborativo e em contínuo aperfeiçoamento}, alinhado aos princípios da
ciência aberta.

Espera-se que sua adoção contribua para:

\begin{itemize}
\tightlist
\item
  Ampliar a \textbf{produção científica} baseada nos dados do RHC
\item
  Fortalecer a \textbf{vigilância oncológica} no país
\item
  Subsidiar a formulação de \textbf{políticas de saúde} mais bem
  fundamentadas em evidências epidemiológicas robustas
\end{itemize}

\begin{center}\rule{0.5\linewidth}{0.5pt}\end{center}

\begin{tcolorbox}[enhanced jigsaw, arc=.35mm, bottomrule=.15mm, bottomtitle=1mm, breakable, colback=white, colbacktitle=quarto-callout-tip-color!10!white, colframe=quarto-callout-tip-color-frame, coltitle=black, left=2mm, leftrule=.75mm, opacityback=0, opacitybacktitle=0.6, rightrule=.15mm, title=\textcolor{quarto-callout-tip-color}{\faLightbulb}\hspace{0.5em}{Como contribuir}, titlerule=0mm, toprule=.15mm, toptitle=1mm]

O RHCgen é um projeto de código aberto. Sugestões, correções e
contribuições são bem-vindas pelo repositório do GitHub:

\begin{quote}
\url{https://github.com/andersonlineu/RHCgen}
\end{quote}

\end{tcolorbox}

\bookmarksetup{startatroot}

\chapter*{Referências}\label{referuxeancias}
\addcontentsline{toc}{chapter}{Referências}

\markboth{Referências}{Referências}

\protect\phantomsection\label{refs}
\begin{CSLReferences}{1}{1}
\bibitem[\citeproctext]{ref-acs2025}
ACS. 2025. \emph{National Cancer Database}.
\url{https://www.facs.org/quality-programs/cancer-programs/national-cancer-database/}.

\bibitem[\citeproctext]{ref-brasil2025}
Brasil. 2025. \emph{Programa Capacitação e Atualização em Abordagens do
Espaço em Análises de Saúde Pública}.
\url{https://www.capacita.geosaude.icict.fiocruz.br/links.php}.

\bibitem[\citeproctext]{ref-bray2024}
{Bray, Freddie, Mathieu Laversanne, Hyuna Sung, et al.} 2024. {``Global
cancer statistics 2022: {GLOBOCAN} estimates of incidence and mortality
worldwide for 36 cancers in 185 countries''}. \emph{CA: A Cancer Journal
for Clinicians} 74: 229--63. \url{https://doi.org/10.3322/caac.21834}.

\bibitem[\citeproctext]{ref-cancertoday2025}
Cancer Today. 2025. \emph{Cancer Today}.
\url{https://gco.iarc.who.int/today/}.

\bibitem[\citeproctext]{ref-chan2018}
Chan, Benjamin K. C. 2018. {``Data Analysis Using {R} Programming''}.
\emph{Advances in Experimental Medicine and Biology} 1082: 47--122.
\url{https://doi.org/10.1007/978-3-319-93791-5_2}.

\bibitem[\citeproctext]{ref-decit2007}
Decit. 2007. {``Integração de informações dos registros de câncer
brasileiros''}. \emph{Revista de Saúde Pública} 41: 865--68.
\url{https://doi.org/10.1590/S0034-89102007000500024}.

\bibitem[\citeproctext]{ref-ecis2025}
ECIS. 2025. \emph{European Cancer Information System}.
\url{https://ecis.jrc.ec.europa.eu/en}.

\bibitem[\citeproctext]{ref-eortc2017}
EORTC. 2017. \emph{European Organisation For Research And Treatment Of
Cancer}. \url{https://www.eortc.org/}.

\bibitem[\citeproctext]{ref-inca2000}
INCA. 2000. \emph{Registro Hospitalar de Câncer: Rotinas e
Procedimentos}. 1º ed. Ministério da Saúde.
\url{https://www.inca.gov.br/publicacoes/manuais/registros-hospitalares-de-cancer}.

\bibitem[\citeproctext]{ref-inca2010}
INCA. 2010. \emph{Registros Hospitalares de Câncer: Planejamento e
Gestão}. 2º ed. Ministério da Saúde.
\url{https://www.inca.gov.br/sites/ufu.sti.inca.local/files//media/document//registros-hospitalares-de-cancer-2010.pdf}.

\bibitem[\citeproctext]{ref-inca2011}
INCA. 2011. \emph{{IntegradorRHC}: ferramenta para a vigilância
hospitalar de câncer no Brasil}. INCA.
\url{https://www.inca.gov.br/sites/ufu.sti.inca.local/files//media/document//folder-integrador-rhc-2011.pdf}.

\bibitem[\citeproctext]{ref-inca2020}
INCA. 2020. \emph{Dicionário das Variáveis da Base de Dados do Sistema
de Informação de Registros Hospitalares de Câncer ({SisRHC})}.
\url{https://irhc.inca.gov.br/RHCNet/documentosTabWin.action?doc=dicionario.dados}.

\bibitem[\citeproctext]{ref-jalal2017}
Jalal, Hawre, Petros Pechlivanoglou, Eline Krijkamp, Fernando
Alarid-Escudero, Eva Enns, e M. G. Myriam Hunink. 2017. {``An Overview
of {R} in Health Decision Sciences''}. \emph{Medical Decision Making}
37: 735--46. \url{https://doi.org/10.1177/0272989X16686559}.

\bibitem[\citeproctext]{ref-jorge2010}
Jorge, Maria Helena Prado de Mello, Ruy Laurenti, e Sabina Léa Davidson
Gotlieb. 2010. {``Avaliação dos sistemas de informação em saúde no
Brasil''}. \emph{Cadernos de Saúde Coletiva}.
\url{https://pesquisa.bvsalud.org/portal/resource/pt/lil-621256}.

\bibitem[\citeproctext]{ref-lowndes2017}
{Lowndes, Julia S. Stewart, Benjamin D. Best, Courtney Scarborough, et
al.} 2017. {``Our path to better science in less time using open data
science tools''}. \emph{Nature Ecology \& Evolution} 1: 160.
\url{https://doi.org/10.1038/s41559-017-0160}.

\bibitem[\citeproctext]{ref-morais2021}
Morais, João Henrique de Araujo, Tulio César Ramos de Oliveira
Konstantyner, Álvaro Luiz Fazenda, Alessandro Sala, e Cleverson Bernardo
Martins. 2021. {``Rtabnetsp: pacote {R} para extração de indicadores de
saúde do estado de {São Paulo}''}. \emph{Epidemiologia e Serviços de
Saúde} 30: e2020576. \url{https://doi.org/10.1590/S1679-4974202100020}.

\bibitem[\citeproctext]{ref-pinto2012}
Pinto, Iêda Vieira, Dandara Novaes Ramos, Maria do Carmo Esteves da
Costa, Cristiane Bárbara Torres Ferreira, e Mariléa Santos Rebelo. 2012.
{``Completude e consistência dos dados dos registros hospitalares de
câncer no Brasil''}. \emph{Cadernos de Saúde Coletiva}, 113--20.

\bibitem[\citeproctext]{ref-rcoreteam2024}
R Core Team. 2024. \emph{R: A Language and Environment for Statistical
Computing}. R Foundation for Statistical Computing.
\url{https://www.r-project.org/}.

\bibitem[\citeproctext]{ref-romero2006}
Romero, Dalia Elena, e Cynthia Braga da Cunha. 2006. {``Avaliação da
qualidade das variáveis sócio-econômicas e demográficas dos óbitos de
crianças menores de um ano registrados no Sistema de Informações sobre
Mortalidade do Brasil (1996/2001)''}. \emph{Cadernos de Saúde Pública}
22: 673--81. \url{https://doi.org/10.1590/S0102-311X2006000300022}.

\bibitem[\citeproctext]{ref-rstudio2024}
RStudio. 2024. \emph{{RStudio}: Integrated Development Environment for
{R}}. Posit Software, PBC. \url{https://www.r-project.org/}.

\bibitem[\citeproctext]{ref-saldanha2019}
Saldanha, Raphael de Freitas, Ronaldo Rocha Bastos, e Christovam
Barcellos. 2019. {``Microdatasus: pacote para download e
pré-processamento de microdados do Departamento de Informática do {SUS}
({DATASUS})''}. \emph{Cadernos de Saúde Pública} 35: e00032419.
\url{https://doi.org/10.1590/0102-311X00032419}.

\bibitem[\citeproctext]{ref-seer2025}
SEER. 2025. \emph{{SEER} Program}.
\url{https://seer.cancer.gov/about/overview.html}.

\bibitem[\citeproctext]{ref-silva2009}
Silva, Natalino Pires da. 2009. {``A utilização dos programas {TABWIN} e
{TABNET} como ferramentas de apoio à disseminação das informações em
saúde''}. Fundação Oswaldo Cruz.
\url{https://www.arca.fiocruz.br/handle/icict/2300}.

\end{CSLReferences}


\backmatter


\end{document}
